% ch5.tex -- Module 6, Chapter 5: Applications of the motor effect
\chapter[Motors, Generators and Braking]{Applications of the Motor Effect: Motors, Generators and Braking}\label{ch5:chap}
\covers{M6.IQ4.01, M6.IQ4.02, M6.IQ4.03}

% ---- chapter-local TikZ helpers (prefix chfive) ----
\tikzset{
  chfiveIn/.pic={\draw[hscink, line width=0.7pt, fill=white] (0,0) circle (0.16);
    \draw[hscink, line width=0.7pt] (-0.11,-0.11) -- (0.11,0.11) (-0.11,0.11) -- (0.11,-0.11);},
  chfiveOut/.pic={\draw[hscink, line width=0.7pt, fill=white] (0,0) circle (0.16);
    \fill[hscink] (0,0) circle (0.05);},
  chfiveWire/.style={hscink, line width=1.4pt},
  pics/chfivePoles/.style={code={
    \fill[hscforce!15, draw=hscink, line width=0.5pt] (-2.05,-1.15) rectangle (-1.8,1.15);
    \node[font=\small\bfseries] at (-1.925,1.38) {N};
    \fill[hscvel!15, draw=hscink, line width=0.5pt] (1.8,-1.15) rectangle (2.05,1.15);
    \node[font=\small\bfseries] at (1.925,1.38) {S};
    \foreach \yy in {-0.95,-0.35,0.35,0.95}
      \draw[field, hscfield!45] (-1.8,\yy) -- (1.8,\yy);}},
}

\inquiry{How has knowledge about the motor effect been applied to technological advances?}

\begin{goals}
  \item calculate the torque on a coil and give the job of each part of a DC motor
  \item explain back emf, the start-up current and the maximum speed of a motor
  \item compare DC and AC generators and sketch their output graphs
  \item explain how AC induction motors and magnetic brakes work
  \item use Lenz's law and conservation of energy for all of these devices
\end{goals}

%% ------------------------------------------------------------------
\section{Torque on a coil and the simple DC motor}\label{ch5:sec:torque}
\covers{M6.IQ4.01}

\begin{bigidea}
A coil carrying a current in a magnetic field feels opposite forces on its two
sides, so it turns. The torque is largest when the plane of the coil is parallel to
the field. In a DC motor, a split-ring commutator reverses the current every half
turn, so the coil keeps turning the same way.
\end{bigidea}

\begin{keywords}
  \kw{Normal}{a line at \ang{90} to the flat face (the plane) of the coil.}
  \kw{Torque ($\tau$)}{the turning effect of a force. Unit: newton metre (\si{N.m}).}
  \kw{Stator}{the part of a motor that does not move. It makes the magnetic field.}
  \kw{Rotor (armature)}{the part that turns: a coil on an iron core, fixed to the axle.}
  \kw{Split-ring commutator}{two half-rings that reverse the current every half turn.}
  \kw{Brushes}{fixed contacts that connect the spinning coil to the supply.}
\end{keywords}

Every electric motor, from a phone's vibration motor to a train, uses one idea. A
current-carrying wire in a magnetic field feels a force. Wind the wire into a
\term{coil}: many loops of wire, where each loop is one \emph{turn}. Put the coil on an
axle, and the forces make it turn.

\begin{refresher}{Force on a wire, and torque}
\begin{itemize}
  \item A wire of length $l$ with current $I$ in a field $B$ feels a force
        $F = lIB\sin\theta$ (\cref{ch2:chap}). The right-hand palm rule gives its
        direction: fingers along $B$, thumb along $I$, and the palm pushes along $F$.
  \item Torque: $\tau = r_\perp F = rF\sin\theta$ (on the sheet). $r_\perp$ is the
        \term{lever arm}: the perpendicular distance from the axis to the line of the
        force. (Here $\theta$ is the angle between $r$ and $F$. In the coil formula
        below, $\theta$ means a different angle.)
\end{itemize}
\end{refresher}

\subsection{Why a coil turns}
Put a rectangular coil between the poles of a magnet, on an \term{axle} (a rod through
its middle). \Cref{ch5:fig:endon} shows the coil end-on: you look along the axle.

\begin{figure}
  \centering
  \begin{tikzpicture}[x=1cm, y=1cm]
    % (a) theta = 90
    \begin{scope}
      \pic {chfivePoles};
      \draw[guide, ->] (1.06,0) -- (1.62,0) node[pos=0.45, below, font=\footnotesize, text=hscmuted] {$B$};
      \draw[->, hscmuted, line width=0.7pt, dashed] (0,0) -- (90:1.25) node[right, font=\footnotesize, text=hscmuted] {normal};
      \draw[hscmuted] (0.42,0) arc (0:90:0.42);
      \node[font=\footnotesize, text=hscmuted] at (45:0.62) {$\theta$};
      \draw[hscink, line width=2.2pt] (180:0.85) -- (0:0.85);
      \pic at (180:0.85) {chfiveOut};
      \pic at (0:0.85) {chfiveIn};
      \draw[force] ($(180:0.85)+(0,0.18)$) -- ++(0,0.62) node[left, text=hscforce, font=\footnotesize] {$F$};
      \draw[force] ($(0:0.85)+(0,-0.18)$) -- ++(0,-0.62) node[right, text=hscforce, font=\footnotesize] {$F$};
      \fill[hscink] (0,0) circle (1.6pt);
      \node[font=\footnotesize, align=center, anchor=north] at (0,-1.7)
        {(a) Plane parallel to $B$\\($\theta = \ang{90}$): largest lever\\arm, largest torque.};
    \end{scope}
    % (b) theta = 50
    \begin{scope}[xshift=4.75cm]
      \pic {chfivePoles};
      \draw[guide, ->] (0,0) -- (1.3,0) node[pos=0.85, below, font=\footnotesize, text=hscmuted] {$B$};
      \draw[->, hscmuted, line width=0.7pt, dashed] (0,0) -- (50:1.25) node[above, font=\footnotesize, text=hscmuted] {normal};
      \draw[hscmuted] (0.42,0) arc (0:50:0.42);
      \node[font=\footnotesize, text=hscmuted] at (25:0.62) {$\theta$};
      \draw[hscink, line width=2.2pt] (140:0.85) -- (-40:0.85);
      \draw[guide] (140:0.85) -- ({0.85*cos(140)},0);
      \draw[hscacc, line width=1.4pt] ({0.85*cos(140)},0) -- (0,0);
      \node[font=\footnotesize, text=hscacc, below] at ({0.425*cos(140)},0) {$r_\perp$};
      \pic at (140:0.85) {chfiveOut};
      \pic at (-40:0.85) {chfiveIn};
      \draw[force] ($(140:0.85)+(0,0.18)$) -- ++(0,0.62) node[left, text=hscforce, font=\footnotesize] {$F$};
      \draw[force] ($(-40:0.85)+(0,-0.18)$) -- ++(0,-0.62) node[right, text=hscforce, font=\footnotesize] {$F$};
      \fill[hscink] (0,0) circle (1.6pt);
      \node[font=\footnotesize, align=center, anchor=north] at (0,-1.7)
        {(b) In between\\($\theta = \ang{50}$): smaller lever\\arm, smaller torque.};
    \end{scope}
    % (c) theta = 0
    \begin{scope}[xshift=9.5cm]
      \pic {chfivePoles};
      \draw[->, hscmuted, line width=0.7pt, dashed] (0,0) -- (0:1.25) node[pos=0.62, above, font=\scriptsize, text=hscmuted] {normal};
      \draw[hscink, line width=2.2pt] (90:0.85) -- (-90:0.85);
      \pic at (90:0.85) {chfiveOut};
      \pic at (-90:0.85) {chfiveIn};
      \draw[force] ($(90:0.85)+(0,0.18)$) -- ++(0,0.42) node[right, text=hscforce, font=\footnotesize] {$F$};
      \draw[force] ($(-90:0.85)+(0,-0.18)$) -- ++(0,-0.42) node[right, text=hscforce, font=\footnotesize] {$F$};
      \fill[hscink] (0,0) circle (1.6pt);
      \node[font=\footnotesize, align=center, anchor=north] at (0,-1.7)
        {(c) Plane perpendicular to $B$\\($\theta = 0$): forces in line,\\zero torque.};
    \end{scope}
  \end{tikzpicture}
  \caption{A coil seen end-on, looking along the axle (black dot). $\odot$ = current out
    of the page, $\otimes$ = current into the page. The forces turn the coil clockwise in
    (a) and (b). $\theta$ is the angle between $B$ and the normal.}\label{ch5:fig:endon}
\end{figure}

\begin{itemize}
  \item The two long sides are at \ang{90} to the field. Each turn of wire feels a
        force $lIB$, so with $n$ turns each long side feels $F = nlIB$. Their currents
        flow opposite ways, so one force points up and the other points down.
  \item The two forces are not along the same line. They form a \term{couple}: two
        equal and opposite forces that turn an object without moving it sideways.
  \item The forces on the short end sides point along the axle. They cancel and give
        no torque.
\end{itemize}
The size of each force stays the same as the coil turns. What changes is the lever
arm $r_\perp$. It is largest in (a). It is zero in (c), where both forces act along one
line through the axle. In the formula, $\theta$ is measured from the \emph{normal},
not from the plane of the coil.

\begin{toolkit}{Areas in square metres, and angles}
\begin{itemize}
  \item $\qty{1}{cm^2} = (\qty{e-2}{m})^2 = \qty{e-4}{m^2}$ and
        $\qty{1}{mm^2} = \qty{e-6}{m^2}$. Example: $\qty{12}{cm^2} = 12 \times
        \qty{e-4}{m^2} = \qty{1.2e-3}{m^2}$. Circle: $A = \pi r^2$ (radius, not diameter).
  \item If the field makes an angle $\phi$ with the \emph{plane} of the coil, then
        $\theta = \ang{90} - \phi$. Also $\sin(\ang{90} - \phi) = \cos\phi$ and
        $\cos(\ang{90} - \phi) = \sin\phi$. Example: plane at \ang{35} to $B$ gives
        $\theta = \ang{55}$, and $\sin\ang{55} = \cos\ang{35} = 0.819$.
\end{itemize}
\end{toolkit}

\textbf{Where the formula comes from.} Let each long side have length $l$, and let the
coil have width $w$. Then the area of one turn is $A = lw$. Each long side is
$\tfrac{w}{2}$ from the axle. The forces are vertical, so the lever arm is the
horizontal distance from the axle to the side. The plane of the coil makes an angle
$(\ang{90} - \theta)$ with the field, so
$r_\perp = \tfrac{w}{2}\cos(\ang{90} - \theta) = \tfrac{w}{2}\sin\theta$
(\cref{ch5:fig:endon}(b)).
\begin{align*}
  \tau_\mathrm{one\ side} &= F r_\perp = nlIB \times \tfrac{w}{2}\sin\theta
    && \text{torque from one long side}\\
  \tau &= 2 \times nlIB \times \tfrac{w}{2}\sin\theta
    && \text{both sides turn the coil the same way}\\
  &= nIB\,(lw)\sin\theta
    && \text{the 2 cancels the $\tfrac{1}{2}$; group $l$ and $w$}\\
  \tau &= nIAB\sin\theta && \text{replace $lw$ with $A$}
\end{align*}

\begin{formulacard}{Torque on a coil}
  \cardformula{\tau = nIA_\perp B = nIAB\sin\theta}
  \begin{symbols}
    \sym{$\tau$}{torque on the coil}{\si{N.m}}
    \sym{$n$}{number of turns in the coil}{no unit}
    \sym{$I$}{current in the coil}{\si{A}}
    \sym{$A$}{area of one turn of the coil}{\si{m^2}}
    \sym{$A_\perp$}{$A\sin\theta$. It equals $A$ when the plane of the coil is
      parallel to $B$, and 0 when the plane is perpendicular to $B$}{\si{m^2}}
    \sym{$B$}{magnetic field strength}{\si{T}}
    \sym{$\theta$}{angle between $B$ and the normal to the coil}{degree (\si{\degree})}
  \end{symbols}
  \inwords{more turns, more current, a bigger coil or a stronger field give more
    torque, and the torque depends on how the coil faces the field.}
  \whentouse{a flat coil of any shape that turns on an axle in a uniform field. In a
    radial field (see \emph{Radial fields} below), $\sin\theta = 1$ in every position.}
  \whennottouse{only one side of a loop is inside the field. Then find $F = lIB$ for
    that side and use $\tau = r_\perp F$.}
  \rearranged{$\tau_\mathrm{max} = nIAB$ (plane $\parallel B$) \qquad
    $I = \dfrac{\tau}{nAB\sin\theta}$ \qquad $B = \dfrac{\tau}{nIA\sin\theta}$}
  \tensecond{$n = 10$, $I = \qty{2.0}{A}$, $A = \qty{0.010}{m^2}$, $B = \qty{0.50}{T}$,
    plane parallel to $B$: $\tau = 10 \times 2.0 \times 0.010 \times 0.50 \times 1
    = \qty{0.10}{N.m}$.}
  \trap{$\theta$ is measured from the \emph{normal}. If the plane of the coil makes
    \ang{30} with the field, use $\theta = \ang{60}$. Change \si{cm^2} to \si{m^2}
    (multiply by \num{e-4}). Put your calculator in degree mode.}
\end{formulacard}

\begin{workedexample}{Maximum torque on a motor coil}\label{ch5:we:maxtorque}
A motor coil has 150 turns. Each turn is a rectangle \qty{4.0}{cm} long and
\qty{2.5}{cm} wide. The \qty{4.0}{cm} sides are parallel to the axle. The coil carries
\qty{0.80}{A} in a uniform \qty{0.25}{T} field. Calculate the maximum torque on the coil.
\given{$n = 150$, \ $I = \qty{0.80}{A}$, \ $B = \qty{0.25}{T}$.
  Area: $A = \qty{4.0}{cm} \times \qty{2.5}{cm} = \qty{10}{cm^2}
  = 10 \times \qty{e-4}{m^2} = \qty{1.0e-3}{m^2}$.
  Maximum torque: plane parallel to $B$, so $\theta = \ang{90}$.}
\find{the maximum torque $\tau_\mathrm{max}$}
\formula{$\tau = nIAB\sin\theta$}
\substitute{$\tau = 150 \times \qty{0.80}{A} \times \qty{1.0e-3}{m^2} \times
  \qty{0.25}{T} \times \sin\ang{90} = \qty{0.030}{N.m}$}
\answer{\final{$\tau_\mathrm{max} = \qty{3.0e-2}{N.m}$} (2 significant figures)}
\checkstep{each \qty{4.0}{cm} side feels $F = nlIB = 150 \times 0.040 \times 0.80
  \times 0.25 = \qty{1.2}{N}$. Each side is half the width from the axle:
  $r = 0.025 \div 2 = \qty{0.0125}{m}$. Two sides:
  $\tau = 2 \times 1.2 \times 0.0125 = \qty{0.030}{N.m}$, the same answer \cmark.}
\end{workedexample}

\begin{figure}
  \centering
  \begin{tikzpicture}[x=1cm, y=1cm]
    \pic {chfivePoles};
    \draw[hscink, line width=2.2pt] (180:0.85) -- (0:0.85);
    \pic at (180:0.85) {chfiveIn};
    \pic at (0:0.85) {chfiveOut};
    \fill[hscink] (0,0) circle (1.6pt);
    \node[font=\scriptsize, text=hscink] at (-0.85,-0.62) {left side};
    \node[font=\scriptsize, text=hscink] at (0.85,-0.62) {right side};
    \node[font=\footnotesize, text=hscfield] at (1.55,0.6) {$B$};
  \end{tikzpicture}
  \caption{The coil in \cref{ch5:we:whichway}, seen end-on. $\otimes$ = current into the
    page, $\odot$ = current out of the page.}\label{ch5:fig:whichway}
\end{figure}

\begin{workedexample}{Which way does the coil turn?}\label{ch5:we:whichway}
\Cref{ch5:fig:whichway} shows a motor coil seen end-on, from the commutator end. The
field points from N (left) to S (right). The current flows into the page in the
left-hand side and out of the page in the right-hand side. Find the direction of the
force on each long side, and the direction in which the coil turns.
\given{$B$ points to the right. Left side: current into the page. Right side: current
  out of the page.}
\find{the force on each side, then the direction of turning}
\formula{the right-hand palm rule: fingers along $B$, thumb along $I$, and the palm
  pushes in the direction of $F$}
\substitute{Left side: fingers to the right, thumb into the page. The palm faces down,
  so the force on the left side points \emph{down}.\par
  Right side: fingers to the right, thumb out of the page. The palm faces up, so the
  force on the right side points \emph{up}.\par
  The left side is pushed down and the right side is pushed up, so the coil turns
  anticlockwise.}
\answer{left side \final{down}, right side \final{up}; the coil turns
  \final{anticlockwise}, seen from the commutator end.}
\checkstep{in \cref{ch5:fig:endon}(a) the currents flow the other way round, and that
  coil turns clockwise. Reversing the current reverses both forces \cmark.}
\end{workedexample}

\subsection{The simple DC motor}
\textbf{A coil on its own will not keep spinning.} Start the coil in position (a) of
\cref{ch5:fig:endon}. The torque turns it to position (c), where the torque is zero.
The coil is moving, so it swings a little past (c). Now the forces turn it
\emph{back}. So the coil rocks back and forth, then stops in position (c). To keep it
turning one way, the current must reverse each time the coil passes position (c).

\Cref{ch5:fig:motor} shows a simple \term{DC motor}: a device that changes electrical
energy from a DC supply into rotational kinetic energy. \Cref{ch5:tab:parts} gives the
job of each part.

\begin{figure}
  \centering
  \begin{tikzpicture}[x={(1cm,0cm)}, y={(0cm,1cm)}, z={(0.42cm,0.30cm)}]
    \def\mz{2.8}
    % N magnet (left)
    \filldraw[draw=hscink, line width=0.5pt, fill=hscforce!28] (-2.6,-1,0) -- (-2.6,-1,\mz) -- (-2.6,1,\mz) -- (-2.6,1,0) -- cycle;
    \filldraw[draw=hscink, line width=0.5pt, fill=hscforce!14] (-3.7,1,0) -- (-2.6,1,0) -- (-2.6,1,\mz) -- (-3.7,1,\mz) -- cycle;
    \filldraw[draw=hscink, line width=0.5pt, fill=hscforce!22] (-3.7,-1,0) -- (-2.6,-1,0) -- (-2.6,1,0) -- (-3.7,1,0) -- cycle;
    \node[font=\bfseries] at (-3.15,0,0) {N};
    % field lines
    \draw[field] (-2.6,0.62,1.4) -- (2.6,0.62,1.4);
    \draw[field] (-2.6,-0.92,0.25) -- (2.6,-0.92,0.25); % low and near the front: misses the commutator
    % S magnet (right)
    \filldraw[draw=hscink, line width=0.5pt, fill=hscvel!14] (2.6,1,0) -- (3.7,1,0) -- (3.7,1,\mz) -- (2.6,1,\mz) -- cycle;
    \filldraw[draw=hscink, line width=0.5pt, fill=hscvel!28] (3.7,-1,0) -- (3.7,-1,\mz) -- (3.7,1,\mz) -- (3.7,1,0) -- cycle;
    \filldraw[draw=hscink, line width=0.5pt, fill=hscvel!22] (2.6,-1,0) -- (3.7,-1,0) -- (3.7,1,0) -- (2.6,1,0) -- cycle;
    \node[font=\bfseries] at (3.15,0,0) {S};
    % axle
    \draw[hscmuted, line width=0.9pt, dashed] (0,0,-1.9) -- (0,0,3.4);
    % coil: current enters the right side (into the page), leaves by the left side
    \draw[chfiveWire, midarrow] (1.4,0,0.3) -- (1.4,0,2.3);
    \draw[chfiveWire, midarrow] (1.4,0,2.3) -- (-1.4,0,2.3);
    \draw[chfiveWire, midarrow] (-1.4,0,2.3) -- (-1.4,0,0.3);
    \draw[chfiveWire] (1.4,0,0.3) -- (0.12,0,0.3) -- (0.12,0,-0.9);
    \draw[chfiveWire] (-1.4,0,0.3) -- (-0.12,0,0.3) -- (-0.12,0,-0.9);
    \node[font=\footnotesize, text=hscink] at (1.75,-0.2,1.6) {$I$};
    \node[font=\footnotesize, text=hscink] at (-1.62,0.25,1.6) {$I$};
    % forces
    \draw[force] (1.4,-0.08,1.3) -- (1.4,-1.0,1.3) node[right, text=hscforce, font=\footnotesize] {$F$};
    \draw[force] (-1.4,0.08,1.3) -- (-1.4,1.0,1.3) node[left, text=hscforce, font=\footnotesize] {$F$};
    \node[font=\footnotesize, text=hscfield] at (2.0,0.85,1.4) {$B$};
    % commutator and brushes at z = -0.9
    \begin{scope}[shift={(0,0,-0.9)}]
      \fill[white] (0,0) circle (0.36);
      \draw[hscseriesB, line width=3pt] (100:0.27) arc (100:260:0.27);
      \draw[hscseriesA, line width=3pt] (-80:0.27) arc (-80:80:0.27);
      \filldraw[fill=hscmuted!55, draw=hscink, line width=0.4pt] (-0.78,-0.09) rectangle (-0.34,0.09);
      \filldraw[fill=hscmuted!55, draw=hscink, line width=0.4pt] (0.34,-0.09) rectangle (0.78,0.09);
      \draw[hscink, line width=0.8pt] (-0.78,0) -- (-1.5,0) -- (-1.5,-2.0) -- (-0.36,-2.0);
      \draw[hscink, line width=0.8pt] (0.78,0) -- (1.5,0) -- (1.5,-2.0) -- (0.0,-2.0);
      \draw[hscink, line width=1.6pt] (-0.36,-2.14) -- (-0.36,-1.86);
      \draw[hscink, line width=0.9pt] (0,-2.3) -- (0,-1.7);
      \node[font=\footnotesize] at (0.22,-1.65) {$+$};
      \node[font=\footnotesize] at (-0.58,-1.65) {$-$};
      \node[diagramlabel, anchor=north] at (-0.18,-2.35) {DC supply};
      \draw[hscmuted, line width=0.4pt] (-0.24,-0.22) -- (-1.75,-0.95) node[anchor=north east, diagramlabel, inner sep=1pt, align=right] {split-ring\\commutator};
      \draw[hscmuted, line width=0.4pt] (0.65,-0.09) -- (1.75,-0.95) node[anchor=north west, diagramlabel, inner sep=1pt] {brush};
    \end{scope}
    \draw[hscmuted, line width=0.4pt] (0,0,3.4) -- ++(0.6,0.5) node[anchor=south west, diagramlabel, inner sep=1pt] {axle};
    \draw[hscmuted, line width=0.4pt] (-0.4,0,2.3) -- ++(-0.2,1.05) node[anchor=south, diagramlabel, inner sep=1pt] {coil (rotor)};
    \node[diagramlabel, anchor=south] at (-3.15,1,2.9) {magnet (stator)};
  \end{tikzpicture}
  \caption{A simple DC motor. The current flows into the page along the right side and
    out of the page along the left side. The coil turns clockwise, as seen from the
    commutator end.}\label{ch5:fig:motor}
\end{figure}

\begin{table}
  \caption{The parts of a simple DC motor and their jobs.}\label{ch5:tab:parts}
  \centering
  \small
  \begin{tabular}{@{}L{36mm}L{102mm}@{}}
    \toprule
    \colhead{Part} & \colhead{Its job}\\
    \midrule
    Stator (field magnets) & permanent magnets or electromagnets that do not move. They
      make the magnetic field.\\
    Rotor (armature) & a coil of many turns on a soft-iron core, fixed to the axle. It
      carries the current, and the forces on it make the torque. The iron makes the
      field stronger.\\
    Axle & turns with the rotor and drives the load (a fan, a wheel, a drill bit).\\
    Split-ring commutator & two insulated half-rings, each joined to one end of the
      coil. It reverses the current in the coil every half turn, so the torque always
      turns the coil the same way.\\
    Brushes & fixed carbon (graphite) blocks that press on the commutator. They keep the
      spinning coil connected to the fixed circuit.\\
    DC power supply & a battery or power pack. It gives a current that flows one way.\\
    \bottomrule
  \end{tabular}
\end{table}

\textbf{How the commutator works} (\cref{ch5:fig:commutator}). Each half-ring is joined
to one end of the coil and turns with it. The brushes stay still. Every half turn, each
half-ring moves from one brush to the other. So each end of the coil swaps between the
$+$ and $-$ terminals. The swap happens when the plane of the coil is perpendicular to
$B$. For an instant the brushes touch the gaps, and no current flows. The coil keeps
going because of its \term{rotational inertia}: the tendency of a spinning object to
keep spinning.

\textbf{The dead spot.} A motor with only one coil has a problem. If it stops with its
plane perpendicular to $B$, the torque is zero, so it cannot start by itself. You must
give it a push. With many coils, at least one coil always has a torque, so the motor
starts by itself.

\begin{figure}
  \centering
  \begin{tikzpicture}[x=1cm, y=1cm]
    % (a) start
    \begin{scope}
      \pic {chfivePoles};
      \draw[->, hscmuted, line width=0.8pt] (125:1.45) arc (125:55:1.45);
      \draw[hscink, line width=2.2pt] (180:1.0) -- (0:1.0);
      \fill[white] (0,0) circle (0.36);
      \draw[hscseriesB, line width=2.6pt] (100:0.22) arc (100:260:0.22);
      \draw[hscseriesA, line width=2.6pt] (-80:0.22) arc (-80:80:0.22);
      \filldraw[fill=hscmuted!55, draw=hscink, line width=0.4pt] (-0.5,-0.06) rectangle (-0.26,0.06);
      \filldraw[fill=hscmuted!55, draw=hscink, line width=0.4pt] (0.26,-0.06) rectangle (0.5,0.06);
      \node[font=\scriptsize] at (0.38,0.2) {$+$};
      \node[font=\scriptsize] at (-0.38,0.2) {$-$};
      \pic at (180:1.0) {chfiveOut};
      \draw[hscseriesB, line width=1.2pt] (180:1.0) circle (0.19);
      \pic at (0:1.0) {chfiveIn};
      \draw[hscseriesA, line width=1.2pt] (0:1.0) circle (0.19);
      \node[font=\footnotesize, text=hscseriesB] at (-1.0,-0.42) {X};
      \node[font=\footnotesize, text=hscseriesA] at (1.0,0.42) {Y};
      \draw[force] ($(180:1.0)+(0,0.2)$) -- ++(0,0.6);
      \draw[force] ($(0:1.0)+(0,-0.2)$) -- ++(0,-0.6);
      \node[font=\footnotesize, text width=4.5cm, align=center, anchor=north] at (0,-1.35)
        {(a) X is joined to $-$ through its half-ring: current out of the page in X.};
    \end{scope}
    % (b) switching moment
    \begin{scope}[xshift=4.75cm]
      \pic {chfivePoles};
      \draw[->, hscmuted, line width=0.8pt] (125:1.45) arc (125:55:1.45);
      \draw[hscink, line width=2.2pt] (90:1.0) -- (-90:1.0);
      \fill[white] (0,0) circle (0.36);
      \draw[hscseriesB, line width=2.6pt] (10:0.22) arc (10:170:0.22);
      \draw[hscseriesA, line width=2.6pt] (190:0.22) arc (190:350:0.22);
      \filldraw[fill=hscmuted!55, draw=hscink, line width=0.4pt] (-0.5,-0.06) rectangle (-0.26,0.06);
      \filldraw[fill=hscmuted!55, draw=hscink, line width=0.4pt] (0.26,-0.06) rectangle (0.5,0.06);
      \filldraw[fill=white, draw=hscseriesB, line width=1.2pt] (90:1.0) circle (0.19);
      \filldraw[fill=white, draw=hscseriesA, line width=1.2pt] (-90:1.0) circle (0.19);
      \node[font=\footnotesize, text=hscseriesB] at (0.42,1.0) {X};
      \node[font=\footnotesize, text=hscseriesA] at (0.42,-1.0) {Y};
      \node[font=\footnotesize, text width=4.5cm, align=center, anchor=north] at (0,-1.35)
        {(b) The gaps pass the brushes: no current, no torque. The coil keeps moving.};
    \end{scope}
    % (c) half a turn later
    \begin{scope}[xshift=9.5cm]
      \pic {chfivePoles};
      \draw[->, hscmuted, line width=0.8pt] (125:1.45) arc (125:55:1.45);
      \draw[hscink, line width=2.2pt] (180:1.0) -- (0:1.0);
      \fill[white] (0,0) circle (0.36);
      \draw[hscseriesB, line width=2.6pt] (-80:0.22) arc (-80:80:0.22);
      \draw[hscseriesA, line width=2.6pt] (100:0.22) arc (100:260:0.22);
      \filldraw[fill=hscmuted!55, draw=hscink, line width=0.4pt] (-0.5,-0.06) rectangle (-0.26,0.06);
      \filldraw[fill=hscmuted!55, draw=hscink, line width=0.4pt] (0.26,-0.06) rectangle (0.5,0.06);
      \node[font=\scriptsize] at (0.38,0.2) {$+$};
      \node[font=\scriptsize] at (-0.38,0.2) {$-$};
      \pic at (0:1.0) {chfiveIn};
      \draw[hscseriesB, line width=1.2pt] (0:1.0) circle (0.19);
      \pic at (180:1.0) {chfiveOut};
      \draw[hscseriesA, line width=1.2pt] (180:1.0) circle (0.19);
      \node[font=\footnotesize, text=hscseriesB] at (1.0,0.42) {X};
      \node[font=\footnotesize, text=hscseriesA] at (-1.0,-0.42) {Y};
      \draw[force] ($(180:1.0)+(0,0.2)$) -- ++(0,0.6);
      \draw[force] ($(0:1.0)+(0,-0.2)$) -- ++(0,-0.6);
      \node[font=\footnotesize, text width=4.5cm, align=center, anchor=north] at (0,-1.35)
        {(c) X is now joined to $+$, so its current has reversed. The torque is still
         clockwise.};
    \end{scope}
  \end{tikzpicture}
  \caption{The commutator seen end-on (it is in front of the coil). Each half-ring has
    the colour of the coil side it is joined to. The right brush is joined to the $+$
    terminal. Current flows from the $+$ brush, round the coil, to the $-$ brush. The
    commutator is at the front (towards you), so current leaving the coil through X
    flows towards you: out of the page.}\label{ch5:fig:commutator}
\end{figure}

\textbf{Torque during a turn.} \Cref{ch5:fig:torquegraph} shows the torque on one coil
turning at a steady rate. Without a commutator the torque changes direction (a). With a
split ring it always acts the same way, but it drops to zero twice in every turn (b).
Real motors have many coils at different angles. Each coil has its own pair of
commutator segments (pieces of the split ring). Their torques add up to a smooth torque
that is never zero.

\begin{figure}
  \centering
  \begin{tikzpicture}
    \begin{axis}[hsc, name=ta, scale only axis, width=10.4cm, height=1.3cm,
        xmin=0, xmax=720, ymin=-1.3, ymax=1.3, xtick={0,90,...,720}, xticklabels={},
        ytick={-1,0,1}, yticklabels={$-\tau_\mathrm{max}$,0,$\tau_\mathrm{max}$}]
      \addplot[hscseriesA, domain=0:720, samples=181] {cos(x)};
      \node[anchor=west, font=\footnotesize] at (axis cs:735,0.0) {(a) no commutator};
      \pgfplotsinvokeforeach{90,270,450,630}{\draw[guide] (axis cs:#1,-1.3) -- (axis cs:#1,1.3);}
      \node[anchor=south, font=\scriptsize, text=hscmuted] at (axis cs:90,1.3) {plane $\perp B$};
      \node[anchor=south, font=\scriptsize, text=hscmuted] at (axis cs:270,1.3) {plane $\perp B$};
    \end{axis}
    \begin{axis}[hsc, name=tb, at={(ta.south west)}, anchor=north west, yshift=-0.3cm,
        scale only axis, width=10.4cm, height=1.1cm,
        xmin=0, xmax=720, ymin=-0.2, ymax=1.3, xtick={0,90,...,720}, xticklabels={},
        ytick={0,1}, yticklabels={0,$\tau_\mathrm{max}$}]
      \addplot[hscseriesB, dashed, domain=0:720, samples=181] {abs(cos(x))};
      \node[anchor=west, font=\footnotesize, align=left] at (axis cs:735,0.55) {(b) split ring,\\ uniform field};
    \end{axis}
    \begin{axis}[hsc, name=tc, at={(tb.south west)}, anchor=north west, yshift=-0.3cm,
        scale only axis, width=10.4cm, height=1.1cm,
        xmin=0, xmax=720, ymin=-0.2, ymax=1.3, xtick={0,90,...,720},
        ytick={0,1}, yticklabels={0,$\tau_\mathrm{max}$},
        xlabel={angle turned (degrees), starting with the plane of the coil parallel to $B$}]
      \addplot[hscseriesC, domain=0:720, samples=721, line width=1.2pt] {min(1, abs(cos(x))/0.12)};
      \node[anchor=west, font=\footnotesize, align=left] at (axis cs:735,0.55) {(c) split ring,\\ radial field};
    \end{axis}
  \end{tikzpicture}
  \caption{Torque on a single coil turned at a steady rate, with a constant current.
    (For (a), imagine the coil is turned by hand.) $\tau_\mathrm{max} = nIAB$ in all
    three graphs.}\label{ch5:fig:torquegraph}
\end{figure}

\begin{sidebyside}[0.40]
  \centering
  \begin{tikzpicture}[x=1cm, y=1cm]
    \filldraw[fill=hscforce!18, draw=hscink, line width=0.5pt]
      (-2.6,0.996) -- (-1.187,0.996) arc (140:220:1.55) -- (-2.6,-0.996) -- cycle;
    \filldraw[fill=hscvel!18, draw=hscink, line width=0.5pt]
      (2.6,0.996) -- (1.187,0.996) arc (40:-40:1.55) -- (2.6,-0.996) -- cycle;
    \node[font=\bfseries] at (-2.1,0) {N};
    \node[font=\bfseries] at (2.1,0) {S};
    \filldraw[fill=hscmuted!22, draw=hscink, line width=0.5pt] (0,0) circle (0.9);
    \node[font=\footnotesize, align=center] at (0,0) {soft-iron\\core};
    \foreach \a in {145,160,175,190,205,215} \draw[field] (\a:1.55) -- (\a:0.9);
    \foreach \a in {-35,-20,-5,10,25,35} \draw[field] (\a:0.9) -- (\a:1.55);
    \pic at (150:1.22) {chfiveOut};
    \pic at (-30:1.22) {chfiveIn};
    \draw[force] ($(150:1.22)+(60:0.18)$) -- ++(60:0.7) node[above, text=hscforce, font=\footnotesize] {$F$};
    \draw[force] ($(-30:1.22)+(240:0.18)$) -- ++(240:0.7) node[below, text=hscforce, font=\footnotesize] {$F$};
    \node[diagramlabel, font=\footnotesize, anchor=north] (pp) at (-1.85,-1.08) {curved pole piece};
    \draw[hscmuted, line width=0.4pt] (-1.62,-1.05) -- (205:1.55);
  \end{tikzpicture}
  \captionof{figure}{A radial field. Each side of the coil always moves at \ang{90} to
    the field lines.}\label{ch5:fig:radial}
\nextside
  \textbf{Radial fields.} Many motors use curved \term{pole pieces} (the shaped ends
  of the magnets) and a soft-iron core in the middle (\cref{ch5:fig:radial}). Together
  they make a \term{radial magnetic field}: a field whose lines point straight out from
  the centre, like the spokes of a wheel. At each side of the coil, the field lies along
  the plane of the coil. So the field is at \ang{90} to the normal: $\theta = \ang{90}$
  in every position. The torque stays at its maximum, $nIAB$, for almost the whole turn
  (\cref{ch5:fig:torquegraph}(c)).
\end{sidebyside}
The maximum torque in a radial field is \emph{not} bigger, but the average torque is. So
the motor runs more smoothly and speeds up more quickly.

\begin{workedexample}{A tilted coil, then a radial field}\label{ch5:we:tilted}
A circular coil of radius \qty{2.0}{cm} has 50 turns and carries \qty{1.5}{A} in a
\qty{0.40}{T} field. (a) The field is uniform and the plane of the coil makes \ang{35}
with it. Calculate the torque. (b) A radial-field motor has a rectangular coil with the
same area, number of turns and current. Its radial field has the same strength.
Calculate its torque.
\given{$n = 50$, \ $r = \qty{2.0}{cm} = 2.0 \times \qty{e-2}{m} = \qty{0.020}{m}$,
  \ $I = \qty{1.5}{A}$, \ $B = \qty{0.40}{T}$.
  Area: $A = \pi r^2 = \pi \times 0.020^2 = \qty{1.257e-3}{m^2}$.
  (a) plane at \ang{35} to $B$, so $\theta = \ang{90} - \ang{35} = \ang{55}$.
  (b) radial field: $\theta = \ang{90}$.}
\find{the torque $\tau$ in each case}
\formula{$\tau = nIAB\sin\theta$}
\substitute{\eqlines{nIAB &= 50 \times \qty{1.5}{A} \times \qty{1.257e-3}{m^2}
  \times \qty{0.40}{T} = \qty{0.0377}{N.m}\\
  \text{(a)}\ \tau &= \qty{0.0377}{N.m} \times \sin\ang{55} = 0.0377 \times 0.819
  = \qty{0.0309}{N.m}\\
  \text{(b)}\ \tau &= \qty{0.0377}{N.m} \times \sin\ang{90} = \qty{0.0377}{N.m}}}
\answer{(a) \final{$\tau = \qty{3.1e-2}{N.m}$} \quad
  (b) \final{$\tau = \qty{3.8e-2}{N.m}$} in every position}
\checkstep{(a) is less than (b) \cmark. Using $\sin\ang{35}$ by mistake gives
  \qty{2.2e-2}{N.m}.}
\end{workedexample}

\begin{toolkit}{Proportional reasoning}
``$y \propto x$'' ($y$ is proportional to $x$) means $y$ doubles when $x$ doubles:
$y_2 = y_1 \times \dfrac{x_2}{x_1}$. If $y \propto ab$, then
$\dfrac{y_2}{y_1} = \dfrac{a_2}{a_1} \times \dfrac{b_2}{b_1}$. Example: $\tau \propto IB$.
Making $I$ 3 times larger and $B$ 2 times larger makes $\tau$ $3 \times 2 = 6$ times larger.
\end{toolkit}

\textbf{Only one side in the field.} Some HSC questions show a loop that turns about a
\term{pivot}: the fixed line it turns around. Only the far side of the loop is in the
field (\cref{ch5:fig:oneside}). Then $\tau = nIAB\sin\theta$ does not apply. Find the
force on that one side with $F = lIB\sin\theta$, then use $\tau = r_\perp F$.

\begin{figure}
  \centering
  \begin{tikzpicture}[x=1cm, y=1cm]
    \fill[hscfield!8] (3.65,-1.0) rectangle (5.35,1.0);
    \foreach \yy in {-0.72,0.72} \draw[->, hscfield, line width=0.8pt] (3.7,\yy) -- (5.3,\yy);
    \node[font=\footnotesize, text=hscfield, anchor=south] at (4.5,1.0) {field $B$};
    \draw[hscink, line width=0.6pt, fill=hscmuted!30] (0,0) -- (-0.22,-0.42) -- (0.22,-0.42) -- cycle;
    \fill[ground] (-0.45,-0.57) rectangle (0.45,-0.42);
    \draw[hscink, line width=0.6pt] (-0.45,-0.42) -- (0.45,-0.42);
    \node[diagramlabel, font=\footnotesize, anchor=east] at (-0.3,0.05) {pivot};
    \draw[chfiveWire] (0,0) -- (4.5,0);
    \fill[hscink] (0,0) circle (2pt);
    \pic at (4.5,0) {chfiveIn};
    \draw[force] (4.5,-0.18) -- (4.5,-1.15) node[below, text=hscforce, font=\footnotesize] {$F$};
    \draw[<->, hscacc, line width=0.7pt] (0,0.32) -- (4.5,0.32)
      node[pos=0.4, above, font=\footnotesize, text=hscacc] {$r_\perp$};
    \node[diagramlabel, font=\footnotesize, anchor=north, align=center] at (2.1,-0.1) {the loop,\\seen edge-on};
    \node[diagramlabel, font=\footnotesize, anchor=west, align=left] at (5.45,0)
      {far side of the loop\\(current into the page)};
  \end{tikzpicture}
  \caption{A loop that turns about a pivot, seen from the side (looking along the pivot).
    Only the far side is in the field. The long sides are parallel to $B$, so they feel
    no force.}\label{ch5:fig:oneside}
\end{figure}

\begin{workedexample}{Only one side in the field}\label{ch5:we:oneside}
A wire loop can turn about a pivot, as in \cref{ch5:fig:oneside}. Only its far side is
inside a \qty{0.080}{T} field, at \ang{90} to the field. That side is \qty{6.0}{cm} long,
carries \qty{2.5}{A} and is \qty{12}{cm} from the pivot. The force on it is at \ang{90} to
the line from the pivot. (a) Calculate the torque. (b) The current is halved. How must
$B$ change to keep the same torque?
\given{$l = \qty{6.0}{cm} = \qty{0.060}{m}$, \ $I = \qty{2.5}{A}$, \ $B = \qty{0.080}{T}$,
  \ $\theta = \ang{90}$, \ $r_\perp = \qty{12}{cm} = \qty{0.12}{m}$}
\find{(a) $\tau$; \ (b) the new $B$}
\formula{$F = lIB\sin\theta$ (one side only), then $\tau = r_\perp F$}
\substitute{\eqlines{F &= \qty{0.060}{m} \times \qty{2.5}{A} \times \qty{0.080}{T}
    \times \sin\ang{90} && \text{force on the far side}\\
  &= \qty{0.012}{N}\\
  \tau &= \qty{0.12}{m} \times \qty{0.012}{N} = \qty{1.44e-3}{N.m}
    && \text{force $\times$ lever arm}\\
  \text{(b)}\ \tau &= r_\perp lIB\text{, so } \tau \propto IB
    && \text{$r_\perp$ and $l$ stay the same}}}
\answer{(a) \final{$\tau = \qty{1.4e-3}{N.m}$} \quad (b) $I \times \tfrac{1}{2}$ needs
  $B \times 2$: \final{double $B$ to \qty{0.16}{T}}}
\checkstep{$\tfrac{1}{2} \times 2 = 1$, so the product $IB$, and the torque, stay the
  same \cmark.}
\end{workedexample}

\begin{mistakes}
\begin{itemize}
  \item \xmark\ Taking $\theta$ from the plane of the coil. \cmark\ $\theta$ is between
        $B$ and the normal. Plane parallel to $B$ means $\theta = \ang{90}$.
  \item \xmark\ Forgetting $n$, or using $\qty{1}{cm^2} = \qty{e-2}{m^2}$. \cmark\ Use $n$, and
        $\qty{1}{cm^2} = \qty{e-4}{m^2}$.
  \item \xmark\ ``The current reverses when the plane of the coil is parallel to the
        field.'' \cmark\ It reverses when the plane is \emph{perpendicular} to the field.
  \item \xmark\ ``A radial field increases the maximum torque.'' \cmark\ The maximum is
        the same. It lasts for longer.
\end{itemize}
\end{mistakes}

\begin{examtip}
\begin{itemize}
  \item ``Calculate the torque'' (2 marks): write $\tau = nIAB\sin\theta$, put $A$ in
        \si{m^2}, and give \si{N.m}. If only one side is in the field, use $F = lIB$,
        then $\tau = r_\perp F$.
  \item ``Explain why a split-ring commutator is needed'' (3--4 marks). Without it,
        the torque reverses each time the plane passes the position perpendicular to
        $B$, so the coil rocks and stops. The split ring reverses the current every half
        turn there, so the torque always acts the same way.
\end{itemize}
\end{examtip}

\begin{quickcheck}
  \item A square coil (side \qty{2.5}{cm}, 200 turns, \qty{0.50}{A}) is in a
        \qty{0.30}{T} field. Find the torque (a) at its maximum, (b) when its plane
        makes \ang{60} with the field.
  \item In \cref{ch5:we:whichway}, the supply is connected the other way round.
        (a) Which way does the coil turn now? (b) The magnets are then swapped as well
        (N on the right). Which way does it turn? Explain each answer.
  \item A wire loop can turn about a pivot, as in \cref{ch5:fig:oneside}. Only one side
        of the loop is in a \qty{0.050}{T} field, at \ang{90} to the field. That side is
        \qty{4.0}{cm} long, carries \qty{3.0}{A} and is \qty{15}{cm} from the pivot. The
        force on it is at \ang{90} to the line from the pivot. (a) Find the torque.
        (b) The current is made 4 times larger. How must $B$ change so that the torque is
        only 2 times larger?
  \item Name the part of a DC motor that (a) reverses the current every half turn,
        (b) keeps the contact with the spinning coil, (c) makes the field.
  \item A coil with no commutator is released with its plane parallel to the field.
        Describe and explain its motion.
\end{quickcheck}

%% ------------------------------------------------------------------
\section{Back emf, maximum speed and energy}\label{ch5:sec:backemf}
\covers{M6.IQ4.01, M6.IQ4.03}

\begin{bigidea}
A spinning motor coil is also a generator. Its induced emf, the back emf, opposes the
supply. Back emf explains the large current at start-up, the motor's maximum speed,
and how the motor turns electrical energy into kinetic energy.
\end{bigidea}

\begin{keywords}
  \kw{Back emf ($\varepsilon_\mathrm{back}$)}{the emf induced in a spinning motor coil.
    It opposes the supply emf.}
  \kw{Net emf}{supply emf minus back emf. It pushes the current through the coil.}
  \kw{Load}{whatever the motor drives (fan blades, a drill bit). It resists the turning.}
\end{keywords}

\begin{refresher}{Flux, Faraday's law and Lenz's law (\cref{ch3:chap})}
Flux through a coil: $\Phi = BA\cos\theta$ (unit \si{Wb}). Faraday's law: a changing
flux induces an emf, $\varepsilon = -N\dfrac{\Delta\Phi}{\Delta t}$. A faster change
gives a bigger emf. Lenz's law: the induced emf opposes the change that produced it
(the minus sign).
\end{refresher}

When a motor coil spins, the flux through it keeps changing. So, by Faraday's law, an
emf is induced in the coil. By Lenz's law, this emf opposes the change that caused it:
the rotation driven by the supply. So it acts against the supply. It is called the
\term{back emf} ($\varepsilon_\mathrm{back}$): the emf induced in a spinning motor coil.
Every motor has one, AC or DC. The faster the coil spins, the faster the flux changes.
So the back emf is proportional to the speed: double the speed, double the back emf.

\begin{refresher}{Ohm's law and power (Year 11, and \cref{ch4:chap})}
$V = IR$: the voltage across a resistor equals current $\times$ resistance. $P = VI$: the
power (energy each second) from a supply. Put $V = IR$ into $P = VI$:
$P = I \times IR = I^2R$. This is the heating power in a resistance $R$.
\end{refresher}

\begin{sidebyside}[0.48]
  \centering
  \begin{tikzpicture}[x=1cm, y=1cm]
    \draw[hscink, line width=0.8pt] (0,0) -- (0,0.88);
    \draw[hscink, line width=0.8pt] (0,1.12) -- (0,2) -- (1.3,2);
    \draw[hscink, line width=0.8pt] (2.5,2) -- (3.6,2) -- (3.6,1.12);
    \draw[hscink, line width=0.8pt] (3.6,0.88) -- (3.6,0) -- (0,0);
    \draw[hscink, line width=1.6pt] (-0.12,0.88) -- (0.12,0.88);
    \draw[hscink, line width=0.9pt] (-0.28,1.12) -- (0.28,1.12);
    \node[font=\scriptsize] at (-0.42,1.25) {$+$};
    \node[diagramlabel, anchor=east, align=right] at (-0.4,0.75) {supply\\$V$};
    \draw[hscink, line width=1.6pt] (3.48,0.88) -- (3.72,0.88);
    \draw[hscink, line width=0.9pt] (3.32,1.12) -- (3.88,1.12);
    \node[font=\scriptsize] at (4.02,1.25) {$+$};
    \node[diagramlabel, anchor=west] at (3.9,0.8) {$\varepsilon_\mathrm{back}$};
    \filldraw[fill=white, draw=hscink, line width=0.8pt] (1.3,1.84) rectangle (2.5,2.16);
    \node[diagramlabel, above] at (1.9,2.16) {$R$};
    \draw[->, hscink, line width=0.8pt] (0.25,2.3) -- (1.05,2.3) node[midway, above, font=\footnotesize] {$I$};
    \draw[hscmuted, dashed, line width=0.6pt] (1.1,0.3) rectangle (4.6,2.75);
    \node[diagramlabel, text=hscmuted, font=\footnotesize] at (2.4,0.55) {the motor coil};
  \end{tikzpicture}
  \captionof{figure}{A model of a running motor. The back emf faces the opposite way
    to the supply.}\label{ch5:fig:circuit}
\nextside
  In \cref{ch5:fig:circuit} the coil acts like a resistor $R$ (its wire) joined to a
  ``battery'' that faces the wrong way.
  \begin{itemize}
    \item The \term{net emf}, $V - \varepsilon_\mathrm{back}$, is what pushes the
          current. So $I = \dfrac{V - \varepsilon_\mathrm{back}}{R}$.
    \item The supply voltage $V$ stays the same. The back emf changes with speed.
  \end{itemize}
\end{sidebyside}

\subsection{From switch-on to maximum speed}
\Cref{ch5:fig:startup} shows what happens after switch-on.
\begin{enumerate}
  \item At switch-on the coil is still, so $\varepsilon_\mathrm{back} = 0$ and
        $I = \dfrac{V}{R}$. The resistance is small, so the current and torque are large.
  \item As the motor speeds up, the back emf grows. The current and torque fall. The
        motor still speeds up, but more slowly.
  \item The motor stops speeding up when its torque just balances friction and the
        load. This is its \term{maximum speed} (the top speed it settles at). In a real
        motor the back emf is then a little less than $V$, and a small current still
        flows. In an ideal motor (no friction, no load), the back emf would almost equal
        $V$. The current would then be almost zero.
\end{enumerate}
You may have seen this at home. When a vacuum cleaner or a fridge starts, the lights
can dim for a moment. The motor draws a large current until it speeds up and its back
emf grows.

\begin{figure}
  \centering
  \begin{tikzpicture}
    \begin{axis}[hsc, name=g1, scale only axis, width=5.1cm, height=2.6cm,
        xmin=0, xmax=5, ymin=0, ymax=1.25, xtick=\empty, ytick={0,1},
        yticklabels={0,$V$}, xlabel={time from switch-on}, ylabel={emf}]
      \addplot[hscmuted, dashed, domain=0:5, samples=2] {1};
      \addplot[hscseriesA, domain=0:5, samples=100] {0.9*(1-exp(-x))};
      \node[anchor=south, font=\footnotesize, text=hscmuted] at (axis cs:3.2,1.0) {supply emf $V$};
      \node[anchor=north west, font=\footnotesize, text=hscseriesA, align=left] at (axis cs:1.75,0.68)
        {back emf $\varepsilon_\mathrm{back}$ (same\\shape as the speed)};
    \end{axis}
    \begin{axis}[hsc, name=g2, at={(g1.south east)}, anchor=south west, xshift=1.9cm,
        scale only axis, width=5.1cm, height=2.6cm,
        xmin=0, xmax=5, ymin=0, ymax=1.25, xtick=\empty, ytick={0,1},
        yticklabels={0,$\dfrac{V}{R}$}, xlabel={time from switch-on}, ylabel={current}]
      \addplot[hscseriesB, domain=0:5, samples=100] {0.1+0.9*exp(-x)};
      \addplot[hscseriesC, densely dotted, line width=1.2pt, domain=0:5, samples=100] {exp(-x)};
      \node[anchor=south west, font=\footnotesize, text=hscseriesB] at (axis cs:2.4,0.2) {real motor};
      \node[anchor=south west, font=\scriptsize, text=hscseriesC, align=left, inner sep=1pt] (idl) at (axis cs:0.12,0.03) {ideal\\motor};
      \draw[hscseriesC, line width=0.4pt] (idl.north east) -- (axis cs:0.95,0.39);
      \node[anchor=west, font=\footnotesize, align=left] at (axis cs:0.35,0.95)
        {largest at switch-on;\\ torque has the same shape};
    \end{axis}
  \end{tikzpicture}
  \caption{A DC motor after switch-on. Left: the back emf rises and levels off just below
    the supply emf. Right: the current falls and levels off. A real motor keeps a small
    current to overcome friction.}\label{ch5:fig:startup}
\end{figure}

\textbf{Loads and stalls.} A bigger \term{load} (whatever the motor drives) slows the
motor. The back emf falls, so the current rises. The torque rises until it balances the
load at a lower speed. If the motor jams, it \term{stalls}: it is switched on but cannot
turn. The back emf is zero, so the current is $\dfrac{V}{R}$, which is large. The coil
overheats. The insulation (the thin coat of enamel or plastic on the wire) can melt,
and the motor ``burns out''. A \term{starting resistor} is a variable resistor in series
with the motor. It limits the start-up current. Its resistance is reduced as the motor
speeds up.

\textbf{Torque and speed.} In a radial field, $\tau = nIAB$, so the torque is
proportional to the current. The back emf is proportional to the speed. So each equal
step up in speed adds the same amount of back emf. It lowers the current
$I = \dfrac{V - \varepsilon_\mathrm{back}}{R}$, and the torque, by the same amount each
time. So a graph of torque against speed is a straight line sloping down. The torque is
largest at zero speed (switch-on).

\subsection{Energy in a DC motor, and why Lenz's law must be true}
\covers{M6.IQ4.03}
The \term{law of conservation of energy} says that energy cannot be made or destroyed;
it only changes form. In a motor, electrical energy from the supply ($P = VI$) becomes
mechanical energy. This is the kinetic energy of the rotor and the work done on the
load. Some energy is wasted. It becomes heat in the coil ($I^2R$). It also becomes heat and sound from
friction in the \term{bearings} (the parts that hold the spinning axle) and the brushes.

Power in is $VI$. The coil wastes $I^2R$ as heat. By conservation of energy, the rest,
$VI - I^2R$, must become mechanical energy. The formula card below shows that this rest
equals $\varepsilon_\mathrm{back}I$. So back emf is not a fault: it is how the motor
does work. A stalled motor has no back emf, so all of its input energy becomes heat.

\textbf{Lenz's law and energy.} Imagine the induced emf \emph{helped} the supply. Then a
faster coil would get a bigger current and a bigger torque. So it would spin faster
still, with no limit. Its kinetic energy and the heat in the coil would both keep
growing. Together they would be more than the energy the supply gives. Energy would be
made from nothing. That breaks the law of conservation of energy. So the induced emf
must oppose the supply: Lenz's law is conservation of energy applied to induction.

\begin{formulacard}[notonsheet]{Current and power in a running motor}
  \cardformula{I = \frac{V - \varepsilon_\mathrm{back}}{R}}
  \begin{symbols}
    \sym{$I$}{current in the motor coil}{\si{A}}
    \sym{$V$}{supply emf (supply voltage)}{\si{V}}
    \sym{$\varepsilon_\mathrm{back}$}{back emf induced in the spinning coil}{\si{V}}
    \sym{$R$}{resistance of the coil}{\si{\ohm}}
  \end{symbols}
  \inwords{the current is set by the net emf (supply minus back emf), not by the supply
    alone.}
  \whentouse{a DC motor that is running, speeding up or slowing down. At switch-on or
    when stalled, $\varepsilon_\mathrm{back} = 0$.}
  \whennottouse{a generator, where the induced emf is the only emf, so
    $I = \dfrac{\varepsilon}{R}$. Never use $I = \dfrac{V}{R}$ for a running motor: it is
    true only at switch-on or when the motor is stalled.}
  \rearranged{$\varepsilon_\mathrm{back} = V - IR$ \qquad power:
    $VI = \varepsilon_\mathrm{back}I + I^2R$ \quad (power in $=$ mechanical power
    $+$ heating power)}
  \tensecond{$V = \qty{6.0}{V}$, $\varepsilon_\mathrm{back} = \qty{4.0}{V}$,
    $R = \qty{2.0}{\ohm}$: $I = \dfrac{6.0 - 4.0}{2.0} = \qty{1.0}{A}$. Power in
    \qty{6.0}{W} $=$ \qty{4.0}{W} mechanical $+$ \qty{2.0}{W} heat.}
  \derive{Ohm's law, $V = IR$ (on the sheet), applies to the coil's resistance. The
    emf that pushes current through $R$ is the net emf.
    \begin{align*}
      V - \varepsilon_\mathrm{back} &= IR && \text{net emf $=$ current $\times$ resistance}\\
      I &= \frac{V - \varepsilon_\mathrm{back}}{R} && \text{divide both sides by $R$}\\
      VI &= \varepsilon_\mathrm{back}I + I^2R && \text{write $V = \varepsilon_\mathrm{back} + IR$, multiply by $I$}
    \end{align*}}
  \trap{the back emf is zero at switch-on, not at top speed. In a motor driven by its
    supply, $\varepsilon_\mathrm{back}$ is less than $V$, so a negative current means a
    mistake. (If an outside force spins it faster, $\varepsilon_\mathrm{back} > V$ and it
    becomes a generator. See regenerative braking in \cref{ch5:sec:induction}.) In $I^2R$
    use the actual current, not $\dfrac{V}{R}$.}
\end{formulacard}

\begin{workedexample}{Start-up, running and energy}\label{ch5:we:startup}
A \qty{12.0}{V} DC motor has a coil resistance of \qty{1.5}{\ohm}.
(a) Calculate the current at switch-on. (b) At full speed it draws \qty{0.80}{A}.
Calculate the back emf. (c) At full speed, find the input power, the heating power and
the mechanical power. Compare the heating with the heating at switch-on.
\given{$V = \qty{12.0}{V}$, \ $R = \qty{1.5}{\ohm}$. At switch-on
  $\varepsilon_\mathrm{back} = 0$. At full speed $I = \qty{0.80}{A}$.}
\find{$I_\mathrm{start}$, \ $\varepsilon_\mathrm{back}$, \ $P_\mathrm{in}$,
  $P_\mathrm{heat}$, $P_\mathrm{mech}$}
\formula{$I = \dfrac{V - \varepsilon_\mathrm{back}}{R}$ \ (so
  $\varepsilon_\mathrm{back} = V - IR$), \quad $VI = \varepsilon_\mathrm{back}I + I^2R$}
\substitute{\eqlines{I_\mathrm{start} &= \frac{\qty{12.0}{V} - 0}{\qty{1.5}{\ohm}}
    = \qty{8.0}{A} && \text{(a) no back emf yet}\\
  \varepsilon_\mathrm{back} &= \qty{12.0}{V} - \qty{0.80}{A} \times \qty{1.5}{\ohm}
    = \qty{10.8}{V} && \text{(b) } \varepsilon_\mathrm{back} = V - IR}\par
  \eqlines{P_\mathrm{in} &= \qty{12.0}{V} \times \qty{0.80}{A} = \qty{9.6}{W}
    && \text{(c) power in, } VI\\
  P_\mathrm{heat} &= (\qty{0.80}{A})^2 \times \qty{1.5}{\ohm} = \qty{0.96}{W}
    && \text{heating in the coil, } I^2R\\
  P_\mathrm{mech} &= \qty{10.8}{V} \times \qty{0.80}{A} = \qty{8.64}{W}
    && \text{the rest, } \varepsilon_\mathrm{back}I\\
  P_\mathrm{heat,\,start} &= (\qty{8.0}{A})^2 \times \qty{1.5}{\ohm} = \qty{96}{W}
    && \text{at switch-on, all heat}}}
\answer{(a) \final{\qty{8.0}{A}} \quad (b) \final{\qty{10.8}{V}} \quad
  (c) \final{\qty{9.6}{W}, \qty{0.96}{W}, \qty{8.6}{W}}. At switch-on, all \qty{96}{W}
  becomes heat: 100 times more heating.}
\checkstep{$8.64 + 0.96 = \qty{9.60}{W}$, so energy is conserved \cmark. The back emf
  is less than \qty{12.0}{V} \cmark.}
\end{workedexample}

\begin{workedexample}{A motor under load}\label{ch5:we:load}
A \qty{24.0}{V} motor with a radial field has a coil resistance of \qty{2.0}{\ohm}.
With no load it spins at \qty{3000}{rpm} (revolutions per minute) with a back emf of
\qty{22.0}{V}.
(a) Calculate the current. (b) A load slows it to \qty{2400}{rpm}. Calculate the new
current. (c) By what factor does the torque change?
\given{$V = \qty{24.0}{V}$, \ $R = \qty{2.0}{\ohm}$, \ $\varepsilon_1 = \qty{22.0}{V}$
  at \qty{3000}{rpm}, then \qty{2400}{rpm}. Back emf $\propto$ speed. Radial field:
  $\tau = nIAB$, so $\tau \propto I$.}
\find{$I_1$, \ $I_2$ and $\dfrac{\tau_2}{\tau_1}$}
\formula{$I = \dfrac{V - \varepsilon_\mathrm{back}}{R}$, \quad
  $\varepsilon_2 = \varepsilon_1 \times \dfrac{\text{speed}_2}{\text{speed}_1}$, \quad
  $\dfrac{\tau_2}{\tau_1} = \dfrac{I_2}{I_1}$}
\substitute{\eqlines{I_1 &= \frac{\qty{24.0}{V} - \qty{22.0}{V}}{\qty{2.0}{\ohm}} = \qty{1.0}{A}
    && \text{(a) net emf $\div$ resistance}\\
  \varepsilon_2 &= \qty{22.0}{V} \times \frac{2400}{3000} = \qty{17.6}{V}
    && \text{back emf $\propto$ speed}\\
  I_2 &= \frac{\qty{24.0}{V} - \qty{17.6}{V}}{\qty{2.0}{\ohm}} = \qty{3.2}{A}
    && \text{(b) new net emf $\div$ resistance}\\
  \frac{\tau_2}{\tau_1} &= \frac{\qty{3.2}{A}}{\qty{1.0}{A}} = 3.2
    && \text{(c) $\tau \propto I$ in a radial field}}}
\answer{(a) \final{\qty{1.0}{A}} \quad (b) \final{\qty{3.2}{A}} \quad
  (c) the torque becomes \final{3.2 times larger}}
\checkstep{slower means less back emf, so more current and more torque to drive the
  load \cmark.}
\end{workedexample}

\begin{mistakes}
\begin{itemize}
  \item \xmark\ ``Back emf comes from the resistance of the coil.'' \cmark\ It comes
        from the changing flux in the spinning coil (Faraday). Lenz gives its direction.
  \item \xmark\ ``The supply voltage falls as the motor speeds up.'' \cmark\ $V$ stays
        the same; the \emph{net} emf falls.
  \item \xmark\ Drawing the current falling to zero for a real motor.
        \cmark\ It levels off above zero, because friction needs some torque.
  \item \xmark\ ``More load means less current.'' \cmark\ More load means less speed,
        less back emf and \emph{more} current.
\end{itemize}
\end{mistakes}

\begin{examtip}
``Explain why a motor can burn out if it jams'' (3 marks) is a chain, one mark per link.
Do not skip the back emf link:
\begin{enumerate}
  \item The spinning coil has a changing flux, so a back emf is induced (Faraday). It
        opposes the supply (Lenz). It is proportional to the speed.
  \item When the motor jams, the back emf falls to zero. So $I = \dfrac{V}{R}$, which is
        large.
  \item The heating power $I^2R$ is very large, so the insulation melts.
\end{enumerate}
For a current--time sketch:
\begin{itemize}
  \item label both axes;
  \item start at $\dfrac{V}{R}$;
  \item draw a curve that levels off above zero;
  \item say if the motor is real or ideal.
\end{itemize}
\end{examtip}

\begin{quickcheck}
  \item A \qty{9.0}{V} motor has a coil resistance of \qty{3.0}{\ohm}. (a) Find the
        current at switch-on. (b) When running, it draws \qty{0.40}{A}. Find the back emf.
  \item Sketch the current in a real DC motor against time from switch-on. Label its
        main features.
  \item A cordless drill has an \qty{18.0}{V} motor with a coil resistance of
        \qty{0.40}{\ohm}. In normal use it draws \qty{3.0}{A}. (a) Find the back emf.
        (b) The bit jams. Find the new current. (c) Compare the heating power in the
        coil in the two cases, and explain why the motor may burn out.
  \item Use the law of conservation of energy to explain why the back emf must oppose
        the supply emf.
\end{quickcheck}

%% ------------------------------------------------------------------
\section{Generators: AC and DC}\label{ch5:sec:generators}
\covers{M6.IQ4.02, M6.IQ4.03}

\begin{bigidea}
A generator is a motor run backwards. An outside force turns a coil in a magnetic
field, and the changing flux induces an emf (Faraday's law). Slip rings give an
alternating output; a split-ring commutator gives a one-way output.
\end{bigidea}

\begin{keywords}
  \kw{Generator}{a device that changes kinetic energy into electrical energy.}
  \kw{Alternating current (AC)}{current that keeps reversing direction.}
  \kw{Direct current (DC)}{current that flows one way only.}
  \kw{Slip rings}{two full rings on the axle. Each is always joined to the same end of
    the coil.}
  \kw{Peak emf}{the largest emf in one cycle.}
\end{keywords}

A \term{generator} is a device that changes kinetic energy into electrical energy. It
has the same parts as a motor: magnets, a coil on an axle, brushes and rings. But the
energy goes the other way. An outside force (a hand, a turbine or an engine) turns the coil.
The output goes to a load such as a lamp (\cref{ch5:fig:acgen}). As the coil turns, the
angle between $B$ and the normal changes. So the flux $\Phi = BA\cos\theta$ changes, and
an emf is induced.

\begin{figure}
  \centering
  \begin{tikzpicture}[x={(1cm,0cm)}, y={(0cm,1cm)}, z={(0.42cm,0.30cm)}, scale=0.95]
    \def\mz{2.8}
    \filldraw[draw=hscink, line width=0.5pt, fill=hscforce!28] (-2.6,-1,0) -- (-2.6,-1,\mz) -- (-2.6,1,\mz) -- (-2.6,1,0) -- cycle;
    \filldraw[draw=hscink, line width=0.5pt, fill=hscforce!14] (-3.7,1,0) -- (-2.6,1,0) -- (-2.6,1,\mz) -- (-3.7,1,\mz) -- cycle;
    \filldraw[draw=hscink, line width=0.5pt, fill=hscforce!22] (-3.7,-1,0) -- (-2.6,-1,0) -- (-2.6,1,0) -- (-3.7,1,0) -- cycle;
    \node[font=\bfseries] at (-3.15,0,0) {N};
    \draw[field] (-2.6,0.62,1.4) -- (2.6,0.62,1.4);
    \draw[field] (-2.6,-0.4,2.6) -- (2.6,-0.4,2.6); % behind the coil: keeps the rings clear
    \filldraw[draw=hscink, line width=0.5pt, fill=hscvel!14] (2.6,1,0) -- (3.7,1,0) -- (3.7,1,\mz) -- (2.6,1,\mz) -- cycle;
    \filldraw[draw=hscink, line width=0.5pt, fill=hscvel!28] (3.7,-1,0) -- (3.7,-1,\mz) -- (3.7,1,\mz) -- (3.7,1,0) -- cycle;
    \filldraw[draw=hscink, line width=0.5pt, fill=hscvel!22] (2.6,-1,0) -- (3.7,-1,0) -- (3.7,1,0) -- (2.6,1,0) -- cycle;
    \node[font=\bfseries] at (3.15,0,0) {S};
    \draw[hscmuted, line width=0.9pt, dashed] (0,0,-2.4) -- (0,0,3.4);
    \draw[hscink, line width=1.6pt] (0,0,3.4) -- (0,0.75,3.4) -- (0,0.75,3.9);
    \draw[->, hscmuted, line width=0.8pt] ([shift={(0,0,3.4)}]150:0.55) arc (150:30:0.55);
    \node[diagramlabel, anchor=west] at (0.5,1.0,3.6) {handle (or turbine)};
    \draw[chfiveWire] (1.4,0,0.3) -- (1.4,0,2.3) -- (-1.4,0,2.3) -- (-1.4,0,0.3) -- (-0.12,0,0.3) -- (-0.12,0,-0.4);
    \draw[chfiveWire] (1.4,0,0.3) -- (0.12,0,0.3) -- (0.12,0,-1.95);
    \draw[hscseriesA, line width=2.6pt] (0,0,-2.0) circle (0.27);
    \draw[hscseriesB, line width=2.6pt] (0,0,-0.45) circle (0.27);
    \filldraw[fill=hscmuted!55, draw=hscink, line width=0.4pt] (-0.07,-0.33,-0.45) rectangle (0.07,-0.62,-0.45);
    \filldraw[fill=hscmuted!55, draw=hscink, line width=0.4pt] (-0.07,-0.33,-2.0) rectangle (0.07,-0.62,-2.0);
    % both leads end on the lamp, in the lamp's plane (z = -1.2)
    \draw[hscink, line width=0.8pt] (0,-0.62,-0.45) -- (0,-1.0,-0.45) -- (1.6,-1.0,-0.45) -- (1.6,-2.05,-0.45) -- (1.6,-2.05,-1.2) -- (0.39,-2.05,-1.2);
    \draw[hscink, line width=0.8pt] (0,-0.62,-2.0) -- (0,-2.05,-2.0) -- (-0.09,-2.05,-1.2);
    \begin{scope}[shift={(0.15,-2.05,-1.2)}]
      \filldraw[fill=hsctip!25, draw=hscink, line width=0.7pt] (0,0) circle (0.24);
      \draw[hscink, line width=0.6pt] (-0.17,-0.17) -- (0.17,0.17) (-0.17,0.17) -- (0.17,-0.17);
    \end{scope}
    \node[diagramlabel, anchor=north west] at (-0.05,-2.35,-1.2) {lamp (load)};
    \node[diagramlabel, anchor=north east, align=right, inner sep=1pt] at (-1.2,-0.8,0) {slip\\ring};
    \draw[hscmuted, line width=0.4pt] (-0.27,0,-2.0) -- (-1.18,-0.85,0);
    \node[diagramlabel, anchor=west, inner sep=1pt] at (0.62,-0.42,0) {slip ring};
    \draw[hscmuted, line width=0.4pt] (0.27,0,-0.45) -- (0.6,-0.38,0);
    \draw[hscmuted, line width=0.4pt] (0.07,-0.5,-0.45) -- (2.75,-1.25,-0.45) node[anchor=west, diagramlabel, inner sep=1pt] {brushes};
    \draw[hscmuted, line width=0.4pt] (-0.4,0,2.3) -- ++(-0.2,1.05) node[anchor=south, diagramlabel, inner sep=1pt] {coil};
  \end{tikzpicture}
  \caption{A simple AC generator. Each slip ring is joined to one end of the coil and
    always touches the same brush.}\label{ch5:fig:acgen}
\end{figure}

\begin{toolkit}{The gradient of a graph is a rate of change}
The gradient (slope) of a graph shows how fast a quantity changes. Where a flux--time
graph is steep, the flux changes quickly, so the emf is large. At the top of a hump the
graph is flat for an instant, so the emf is zero. The emf graph follows the gradient of
the flux graph (with a minus sign).
\end{toolkit}

\Cref{ch5:fig:gengraphs} shows the flux and the emf for a coil turning at a steady
rate. Two positions matter most:
\begin{itemize}
  \item \textbf{Plane perpendicular to $B$:} the flux is largest ($BA$) but not changing
        for an instant (the sides move along the field lines). The emf is \emph{zero}.
  \item \textbf{Plane parallel to $B$:} the flux is zero but changing fastest (the sides
        cut straight across the field lines). The emf is \emph{largest}.
\end{itemize}

\begin{figure}
  \centering
  \begin{tikzpicture}
    \begin{axis}[hsc, name=fa, scale only axis, width=10.2cm, height=1.2cm,
        xmin=0, xmax=2, ymin=-1.3, ymax=1.3, xtick={0,0.25,...,2}, xticklabels={},
        ytick={-1,0,1}, yticklabels={$-BA$,0,$BA$}, ylabel={$\Phi$},
        ylabel style={rotate=-90}]
      \addplot[hscseriesA, domain=0:2, samples=161] {cos(360*x)};
    \end{axis}
    \foreach \f/\ang in {0/90, 0.125/0, 0.25/90, 0.375/0, 0.5/90} {
      \coordinate (ic) at ($(fa.north west)!\f!(fa.north east)+(0,0.95)$);
      \draw[hscink, line width=1.6pt] ($(ic)+({\ang}:0.24)$) -- ($(ic)+({\ang+180}:0.24)$);
      \fill[hscink] (ic) circle (1pt);
    }
    \node[anchor=west, font=\footnotesize, text=hscmuted] at ($(fa.north west)!0.56!(fa.north east)+(0,0.95)$)
      {$\leftarrow$ the coil seen end-on (field $B$ to the right)};
    \node[anchor=south, font=\scriptsize, text=hscmuted, align=center] at ($(fa.north west)!0.0!(fa.north east)+(0,1.22)$) {plane\\$\perp B$};
    \node[anchor=south, font=\scriptsize, text=hscmuted, align=center] at ($(fa.north west)!0.125!(fa.north east)+(0,1.22)$) {plane\\$\parallel B$};
    \begin{axis}[hsc, name=fb, at={(fa.south west)}, anchor=north west, yshift=-0.3cm,
        scale only axis, width=10.2cm, height=1.2cm,
        xmin=0, xmax=2, ymin=-1.3, ymax=1.3, xtick={0,0.25,...,2}, xticklabels={},
        ytick={-1,0,1}, yticklabels={$-\varepsilon_\mathrm{peak}$,0,$\varepsilon_\mathrm{peak}$},
        ylabel={emf}]
      \addplot[hscseriesB, domain=0:2, samples=161] {sin(360*x)};
      \node[anchor=west, font=\footnotesize, align=left] at (axis cs:2.05,0) {AC\\(slip rings)};
    \end{axis}
    \begin{axis}[hsc, name=fc, at={(fb.south west)}, anchor=north west, yshift=-0.3cm,
        scale only axis, width=10.2cm, height=1.2cm,
        xmin=0, xmax=2, ymin=-0.3, ymax=1.3, xtick={0,0.25,...,2},
        xticklabels={0,$T/4$,$T/2$,$3T/4$,$T$,,,,$2T$},
        ytick={0,1}, yticklabels={0,$\varepsilon_\mathrm{peak}$}, ylabel={emf}, xlabel={time}]
      \addplot[hscseriesC, domain=0:2, samples=161] {abs(sin(360*x))};
      \node[anchor=west, font=\footnotesize, align=left] at (axis cs:2.05,0.5) {DC\\(split ring)};
    \end{axis}
  \end{tikzpicture}
  \caption{Flux and output emf for a coil that starts with its plane perpendicular to $B$
    and turns at a steady rate. $T$ is the time for one turn.}\label{ch5:fig:gengraphs}
\end{figure}

\textbf{AC generator.} It has \term{slip rings}: two full rings, each always joined to
the same end of the coil. So each brush always touches the same ring. The emf in the
coil reverses every half turn, so the output reverses too. This is
\term{alternating current (AC)}: current that keeps reversing direction.

\textbf{DC generator.} It looks like the motor in \cref{ch5:fig:motor}. But a load,
such as a lamp, replaces the DC supply, and a handle or turbine turns the axle. It has a
split-ring commutator, as in a motor. The connections
swap every half turn, just when the emf is zero (plane perpendicular to $B$). So the
output never reverses. It is \term{direct current (DC)}: current that flows one way
only. It pulses between zero and the peak; more coils make it smoother. The emf
\emph{inside} the coil still alternates. The split ring only flips the outside
connection.

\textbf{What changes the output?} Turning faster makes the flux change faster, so the
\term{peak emf} (the largest emf in a cycle) is bigger. Each turn also takes less time,
so the frequency is higher. Twice the speed gives twice the peak and half the period.
More turns, a stronger field or a bigger coil raise the peak only. Turning the other
way flips the graph upside down.

\textbf{Lenz's law and energy.} With the circuit open, no current flows, and the coil is
easy to turn. With a lamp connected, a current flows in the coil. This current is in a
magnetic field, so it feels a motor-effect force. By Lenz's law, the force opposes the
rotation, so you must do work against it. That work becomes the electrical energy in
the lamp. A lower-resistance load draws a bigger current, so the handle is harder to
turn. If the force helped the rotation, the generator would make energy from nothing.

\begin{formulacard}{Average emf of a turning coil}
  \cardformula{\varepsilon = -N\frac{\Delta\Phi}{\Delta t}}
  \begin{symbols}
    \sym{$\varepsilon$}{average induced emf}{\si{V}}
    \sym{$N$}{number of turns in the coil (the same as $n$ in the torque formula;
      the sheet uses a capital $N$ here)}{no unit}
    \sym{$\Delta\Phi$}{change in flux through one turn: $\Phi_\mathrm{final} - \Phi_\mathrm{initial}$}{\si{Wb}}
    \sym{$\Delta t$}{time taken for the change}{\si{s}}
  \end{symbols}
  \inwords{the faster the flux through the coil changes, the bigger the emf. The minus
    sign means the emf opposes the change (Lenz's law).}
  \whentouse{a generator coil turns from one position to another. Find each flux with
    $\Phi = BA\cos\theta$ (on the sheet).}
  \whennottouse{to find the emf at one instant (such as the peak). This formula gives
    the \emph{average} emf over $\Delta t$.}
  \rearranged{$|\varepsilon| = \dfrac{N\,|\Delta\Phi|}{\Delta t}$ \qquad
    $\Delta t = \dfrac{N\,|\Delta\Phi|}{|\varepsilon|}$. For a quarter turn between plane
    $\parallel B$ and plane $\perp B$, $|\Delta\Phi| = BA$.}
  \tensecond{$N = 100$, $|\Delta\Phi| = \qty{2.0e-3}{Wb}$ in \qty{0.050}{s}:
    $|\varepsilon| = \dfrac{100 \times 2.0\times10^{-3}}{0.050} = \qty{4.0}{V}$.}
  \trap{``plane parallel to $B$'' means $\theta = \ang{90}$, so $\Phi = 0$, not $BA$.
    Change \si{cm^2} to \si{m^2} first.}
\end{formulacard}

\begin{workedexample}{Average emf for a quarter turn}\label{ch5:we:quarter}
A generator coil has 200 turns, each \qty{5.0}{cm} by \qty{4.0}{cm}, in a \qty{0.15}{T}
field. It turns from the position where its plane is parallel to the field to where its
plane is perpendicular, in \qty{0.010}{s}. Calculate the average emf.
\given{$N = 200$, \ $A = 0.050 \times 0.040 = \qty{2.0e-3}{m^2}$, \ $B = \qty{0.15}{T}$,
  \ $\Delta t = \qty{0.010}{s}$. Start: plane $\parallel B$ ($\theta = \ang{90}$). End:
  plane $\perp B$ ($\theta = 0$).}
\find{the average emf $\varepsilon$}
\formula{$\Phi = BA\cos\theta$, \quad then \quad $\varepsilon = -N\dfrac{\Delta\Phi}{\Delta t}$}
\substitute{\eqlines{\Phi_\mathrm{initial} &= BA\cos\ang{90} = 0\\
  \Phi_\mathrm{final} &= \qty{0.15}{T} \times \qty{2.0e-3}{m^2} \times \cos 0 = \qty{3.0e-4}{Wb}\\
  \varepsilon &= -200 \times \frac{\qty{3.0e-4}{Wb} - 0}{\qty{0.010}{s}} = \qty{-6.0}{V}}}
\answer{the average emf is \final{\qty{6.0}{V}} in size. The minus sign shows that it
  opposes the change in flux (Lenz's law).}
\checkstep{units: $\si{Wb} \div \si{s} = \si{V}$ \cmark. A few volts from a quick
  quarter turn of a hand-sized coil is sensible.}
\end{workedexample}

\begin{toolkit}{Revolutions per minute}
rpm to turns per second: divide by 60 (\qty{3000}{rpm} $= 50$ turns each second). Turns
per second is the frequency $f$ (in \si{Hz}). The time for one turn is $T = \dfrac{1}{f}$
(from $f = \dfrac{1}{T}$ on the sheet).
\end{toolkit}

\begin{workedexample}{Turning faster, and the other way}\label{ch5:we:faster}
An AC generator turns at \qty{1200}{rpm}. Its peak emf is \qty{15}{V}. (a) Calculate the
frequency and period of the output. (b) It is now turned at \qty{2400}{rpm} in the
opposite direction. State the new peak emf and period, and sketch both graphs.
\given{\qty{1200}{rpm}, peak \qty{15}{V}; then \qty{2400}{rpm}, opposite direction}
\find{$f$ and $T$; then the new peak and $T$}
\formula{$f = \dfrac{\text{rpm}}{60}$, \quad $T = \dfrac{1}{f}$, \quad peak emf $\propto$ speed}
\substitute{\eqlines{f &= \frac{1200}{60} = \qty{20}{Hz}, \qquad T = \frac{1}{\qty{20}{Hz}} = \qty{0.050}{s}\\
  f_\mathrm{new} &= \frac{2400}{60} = \qty{40}{Hz}, \qquad T_\mathrm{new} = \frac{1}{\qty{40}{Hz}} = \qty{0.025}{s}\\
  \text{peak}_\mathrm{new} &= \qty{15}{V} \times \frac{2400}{1200} = \qty{30}{V}}}
\answer{(a) \final{\qty{20}{Hz}, \qty{0.050}{s}} \quad (b) \final{\qty{30}{V}, \qty{0.025}{s}},
  and the graph is upside down (\cref{ch5:fig:faster}).}
\checkstep{twice the speed: the flux changes twice as fast (double the emf), and each
  turn takes half the time \cmark.}
\end{workedexample}

\begin{figure}
  \centering
  \begin{tikzpicture}
    \begin{axis}[hsc, scale only axis, width=10.2cm, height=2.2cm,
        xmin=0, xmax=0.1, ymin=-35, ymax=35, xtick={0,0.025,0.05,0.075,0.1},
        xticklabels={0,0.025,0.050,0.075,0.100},
        ytick={-30,-15,0,15,30}, xlabel={time (\si{s})}, ylabel={emf (\si{V})}]
      \addplot[hscseriesA, domain=0:0.1, samples=401] {15*sin(360*20*x)};
      \addplot[hscseriesB, dashed, domain=0:0.1, samples=401] {-30*sin(360*40*x)};
      \node[anchor=west, font=\footnotesize, text=hscseriesA] at (axis cs:0.101,15) {\qty{1200}{rpm}};
      \node[anchor=west, font=\footnotesize, text=hscseriesB, align=left] at (axis cs:0.101,-26) {\qty{2400}{rpm},\\ other way};
    \end{axis}
  \end{tikzpicture}
  \caption{Output of the generator in \cref{ch5:we:faster}.}\label{ch5:fig:faster}
\end{figure}

\textbf{Motors and generators compared.} \Cref{ch5:tab:motorgen} sums up how a DC motor
and a DC generator are alike and how they differ.

\begin{table}
  \caption{DC motor and DC generator compared.}\label{ch5:tab:motorgen}
  \centering
  \small
  \begin{tabular}{@{}L{28mm}L{55mm}L{57mm}@{}}
    \toprule
     & \colhead{DC motor} & \colhead{DC generator}\\
    \midrule
    Same parts & \multicolumn{2}{L{114mm}}{magnets (stator), coil on an axle (rotor),
      split-ring commutator, brushes}\\
    Connected to & a DC supply; drives a load & a turbine or handle; feeds a load\\
    Energy change & electrical $\to$ kinetic & kinetic $\to$ electrical\\
    Main physics & motor effect: force on a current in a field & induction: Faraday's law\\
    Side effect & back emf: it also acts as a generator
      & opposing torque when current flows: it also acts as a motor\\
    \bottomrule
  \end{tabular}
\end{table}

\begin{mistakes}
\begin{itemize}
  \item \xmark\ ``The emf is largest when the flux is largest.'' \cmark\ It is largest
        when the flux is zero but changing fastest (plane parallel to $B$).
  \item \xmark\ ``Turning faster only changes the frequency.'' \cmark\ It raises the
        peak \emph{and} the frequency.
  \item \xmark\ Calling slip rings a commutator. \cmark\ Only the split ring is a
        commutator.
  \item \xmark\ ``A generator is easier to turn when more current is drawn.''
        \cmark\ More current means a bigger opposing force, so it is harder to turn.
\end{itemize}
\end{mistakes}

\begin{examtip}
\textbf{Generator graph sketches} (2--3 marks):
\begin{itemize}
  \item keep the shape: a sine wave for slip rings, one-way humps for a split ring;
  \item multiply the peak by the speed factor, and divide the period by it;
  \item flip the graph if the direction is reversed;
  \item label both axes with units.
\end{itemize}
\textbf{``Explain why the handle is harder to turn with a lamp connected''} (3 marks):
\begin{enumerate}
  \item A current flows in the coil.
  \item The current is in a field, so it feels a force.
  \item By Lenz's law the force opposes the turning.
  \item Your extra work becomes electrical energy (conservation of energy).
\end{enumerate}
\end{examtip}

\begin{quickcheck}[breakable]
  \item State the difference in structure between an AC and a DC generator, and its
        effect on the output.
  \item The flux through a generator coil is zero when its plane is parallel to the
        field. Explain why the emf is largest at this moment.
  \item A generator coil is turned clockwise (seen from the front) in a field that
        points from left to right. At one moment its plane is parallel to the field, so
        its right-hand side moves down. Use Lenz's law and the right-hand palm rule to
        find the direction of the induced current in this side.
  \tcbbreak % layout QA (book build): break here, so "(continued)" is not one lone question
  \item A coil of 50 turns and area \qty{1.2e-3}{m^2} is in a \qty{0.30}{T} field. It
        turns from plane perpendicular to plane parallel to the field in \qty{0.020}{s}.
        Calculate the average emf.
  \item An AC generator turns at \qty{1500}{rpm} and has a peak emf of \qty{8.0}{V}.
        (a) Find the period of its output. (b) It is slowed to \qty{750}{rpm}. State the
        new peak emf and period.
\end{quickcheck}

%% ------------------------------------------------------------------
\section{AC induction motors and magnetic braking}\label{ch5:sec:induction}
\covers{M6.IQ4.02, M6.IQ4.03}

\begin{bigidea}
When a magnetic field and a piece of metal move past each other, currents are induced
in the metal. By Lenz's law, the forces on these currents oppose the \emph{relative}
motion. A moving field drags the metal along (an AC induction motor), and a still field
slows moving metal down (a magnetic brake).
\end{bigidea}

\begin{keywords}
  \kw{Eddy currents}{loops of current induced inside a solid piece of metal when the
    flux through it changes.}
  \kw{Rotating magnetic field}{a field of fixed strength whose direction turns steadily.
    Coils that do not move make it.}
  \kw{Squirrel-cage rotor}{metal bars joined at both ends by end rings, like a cage.}
  \kw{Slip}{the difference between the speed of the field and the speed of the rotor.}
  \kw{Magnetic braking}{slowing a moving metal object with eddy currents, without
    touching it.}
  \kw{Relative motion}{how one thing moves compared with another. A magnet moving past
    still metal, and metal moving past a still magnet, are the same relative motion.}
\end{keywords}

You met eddy currents in \cref{ch3:chap,ch4:chap}. In a transformer they waste energy.
Here they do useful jobs.

A simple demonstration shows the idea. Spin a strong magnet just above a light
aluminium disc that is free to turn. The magnet does not touch the disc, and aluminium
is not magnetic. Yet the disc turns the same way as the magnet, but more slowly. The
moving field changes the flux through parts of the disc, so eddy currents flow
(Faraday). These currents are in the field, so they feel forces. By Lenz's law, the
forces oppose the relative motion, so they drag the disc after the magnet.

\subsection{AC induction motors}
Many motors in fans, pumps, washing machines and factories are \term{AC induction
motors}. In these motors the current in the rotor is induced: no wire supplies it. They
have no brushes and no commutator.

\textbf{Step 1: the stator makes a rotating field.} The stator has three pairs of coils
around a ring (\cref{ch5:fig:stator}). Each pair is joined to its own AC supply.
Together the three supplies are called \term{three-phase AC}: three AC supplies that
peak one after another. They are a third of a cycle (\ang{120}) apart. Each supply is
one \emph{phase}.
\begin{itemize}
  \item At $t_1$ current 1 peaks, so the combined field points along pair 1.
  \item A third of a cycle later ($t_2$), current 2 peaks. The field points along pair 2.
  \item Another third of a cycle later ($t_3$), current 3 peaks. The field points along
        pair 3.
\end{itemize}
The field keeps the same strength while its direction turns. This is a
\term{rotating magnetic field}. No coil moves, but the field really rotates. With one
pair of coils per phase, it makes one full turn in every AC cycle.

\begin{figure}
  \centering
  \begin{tikzpicture}
    \begin{axis}[hsc, name=ph, scale only axis, width=5.4cm, height=2.4cm,
        xmin=0, xmax=1, ymin=-1.25, ymax=1.25, xtick={0,0.5,1},
        xticklabels={0,$\frac{T}{2}$,$T$}, ytick={-1,0,1}, yticklabels={,0,},
        xlabel={time}, ylabel={current}]
      \addplot[hscseriesA, domain=0:1, samples=120] {cos(360*x)};
      \addplot[hscseriesB, dashed, domain=0:1, samples=120] {cos(360*x-120)};
      \addplot[hscseriesC, densely dotted, line width=1.2pt, domain=0:1, samples=120] {cos(360*x-240)};
      \draw[guide] (axis cs:0.3333,-1.25) -- (axis cs:0.3333,1.25);
      \draw[guide] (axis cs:0.6667,-1.25) -- (axis cs:0.6667,1.25);
      \node[font=\scriptsize, anchor=south] at (axis cs:0.01,1.25) {$t_1$};
      \node[font=\scriptsize, anchor=south] at (axis cs:0.3333,1.25) {$t_2$};
      \node[font=\scriptsize, anchor=south] at (axis cs:0.6667,1.25) {$t_3$};
      \node[font=\scriptsize, text=hscseriesA, anchor=south west] at (axis cs:0.03,1.0) {1};
      \node[font=\scriptsize, text=hscseriesB, anchor=south west] at (axis cs:0.35,1.0) {2};
      \node[font=\scriptsize, text=hscseriesC, anchor=south west] at (axis cs:0.6833,1.0) {3};
    \end{axis}
    \begin{scope}[shift={($(ph.east)+(3.6,0.1)$)}, x=1cm, y=1cm]
      \draw[hscmuted!70, line width=4pt] (0,0) circle (1.75);
      \foreach \a/\c/\n in {90/hscseriesA/1, 270/hscseriesA/1, -30/hscseriesB/2, 150/hscseriesB/2, 210/hscseriesC/3, 30/hscseriesC/3} {
        \begin{scope}[rotate=\a]
          \filldraw[fill=\c!20, draw=\c, line width=0.7pt] (1.12,-0.2) rectangle (1.65,0.2);
          \foreach \w in {1.2,1.32,1.44,1.56} \draw[\c, line width=0.7pt] (\w,-0.24) -- (\w,0.24);
        \end{scope}
        \node[font=\scriptsize, text=\c] at (\a:0.95) {\n};
      }
      \draw[->, hscfield!35, line width=2.2pt] (0,0) -- (90:0.7);
      \draw[->, hscfield!65, line width=2.2pt] (0,0) -- (-30:0.7);
      \draw[->, hscfield, line width=2.2pt] (0,0) -- (210:0.7);
      \node[font=\scriptsize, text=hscfield] at (112:0.5) {$t_1$};
      \node[font=\scriptsize, text=hscfield] at (-55:0.45) {$t_2$};
      \node[font=\scriptsize, text=hscfield] at (235:0.5) {$t_3$};
      \draw[->, hscmuted, line width=0.8pt] (115:2.05) arc (115:65:2.05);
      \node[font=\footnotesize, text=hscmuted, anchor=south] at (90:2.08) {field turns};
    \end{scope}
  \end{tikzpicture}
  \caption{Left: the three phase currents. Right: the stator; each coil pair has the
    colour of its phase. The arrows show the combined field at $t_1$, $t_2$ and $t_3$:
    it turns clockwise.}\label{ch5:fig:stator}
\end{figure}

\begin{sidebyside}[0.44]
  \centering
  \begin{tikzpicture}[x=1cm, y=1cm]
    \foreach \a in {0,45,...,315} {
      \pgfmathsetmacro\yy{0.95*sin(\a)}
      \pgfmathsetmacro\xx{0.32*cos(\a)}
      \pgfmathtruncatemacro\isback{cos(\a) < -0.01 ? 1 : 0}
      \ifnum\isback=1
        \draw[hscmuted!60, line width=1.1pt] (\xx,\yy) -- ({\xx+3.0},\yy);
      \fi
    }
    \draw[hscmuted, line width=0.9pt, dashed] (-0.8,0) -- (3.8,0);
    \foreach \a in {0,45,...,315} {
      \pgfmathsetmacro\yy{0.95*sin(\a)}
      \pgfmathsetmacro\xx{0.32*cos(\a)}
      \pgfmathtruncatemacro\isback{cos(\a) < -0.01 ? 1 : 0}
      \ifnum\isback=0
        \draw[hscseriesB!85!black, line width=1.6pt] (\xx,\yy) -- ({\xx+3.0},\yy);
      \fi
    }
    \draw[hscink, line width=2.4pt] (0,0) ellipse (0.32 and 0.95);
    \draw[hscink, line width=2.4pt] (3.0,0) ellipse (0.32 and 0.95);
    \node[diagramlabel, anchor=north] at (0,-1.05) {end ring};
    \node[diagramlabel, anchor=north] at (3.0,-1.05) {end ring};
    \draw[hscmuted, line width=0.4pt] (1.6,0.67) -- (1.9,1.3) node[anchor=south, diagramlabel] {metal bars};
    \node[diagramlabel, anchor=west] at (3.8,0) {axle};
  \end{tikzpicture}
  \captionof{figure}{A squirrel-cage rotor (iron core not shown).}\label{ch5:fig:cage}
\nextside
  \textbf{Step 2: the rotor.} It is a \term{squirrel-cage rotor} (\cref{ch5:fig:cage}):
  copper or aluminium bars joined at both ends by end rings. So currents can flow round
  loops of bars. The bars sit in an iron core made of thin insulated sheets. The sheets
  reduce eddy-current losses in the core, as in a transformer.
\end{sidebyside}

\begin{sidebyside}[0.36]
  \centering
  \begin{tikzpicture}[x=1cm, y=1cm]
    \filldraw[fill=hscforce!15, draw=hscink, line width=0.5pt] (-1.7,1.0) rectangle (1.7,1.3);
    \node[font=\footnotesize\bfseries] at (0,1.15) {N};
    \filldraw[fill=hscvel!15, draw=hscink, line width=0.5pt] (-1.7,-1.3) rectangle (1.7,-1.0);
    \node[font=\footnotesize\bfseries] at (0,-1.15) {S};
    \foreach \xx in {-1.4,-0.8,0.8,1.4} \draw[field] (\xx,1.0) -- (\xx,-1.0);
    \node[font=\footnotesize, text=hscfield] at (-1.1,0.5) {$B$};
    \draw[->, hscmuted, line width=0.9pt] (-0.9,1.55) -- (0.9,1.55)
      node[midway, above, font=\footnotesize, text=hscmuted] {field moves};
    \pic at (0,0) {chfiveOut};
    \draw[force] (0.2,0) -- (0.64,0) node[pos=0.55, above, text=hscforce, font=\footnotesize] {$F$};
    \node[diagramlabel, font=\footnotesize, anchor=north] at (0,-0.22) {rotor bar};
  \end{tikzpicture}
  \captionof{figure}{One rotor bar, end-on. The field sweeps right, past the bar. The
    induced current flows out of the page ($\odot$). The force $F$ on the bar points
    right: the bar is dragged after the field.}\label{ch5:fig:bar}
\nextside
  \textbf{Step 3: why the rotor turns} (\cref{ch5:fig:bar}).
  \begin{enumerate}
    \item The rotating field sweeps past the bars, so the flux through each loop of
          bars keeps changing.
    \item By Faraday's law, emfs are induced in the bars. The end rings complete the
          circuit, so currents flow.
    \item These currents are in the field, so each bar feels a force (the motor
          effect, $F = lIB$).
    \item By Lenz's law, the induced currents oppose the \emph{relative} motion of
          field and rotor. So the forces drag the rotor round in the same direction as
          the field.
  \end{enumerate}
\end{sidebyside}

\textbf{Slip and load.} The rotor always turns a little slower than the field. The
difference in speed is the \term{slip}. If the rotor caught up, there would be no
relative motion. Then there would be no flux change, no current and no force, so
friction would slow the rotor again. A bigger load slows the rotor, so the slip grows.
The flux changes faster, so the induced currents and the torque are larger. The motor
settles at a slightly lower speed.

Slip $=$ speed of the field $-$ speed of the rotor (not on the formula sheet). Example:
the field turns at \qty{1500}{rpm} and the rotor at \qty{1440}{rpm}. Slip
$= 1500 - 1440 = \qty{60}{rpm}$, which is $60 \div 60 = 1$ turn each second.

\textbf{Advantages:} no brushes or commutator, so little wear, no sparking and little
maintenance; tough and cheap. \textbf{Limitations:} the rotor speed stays close to the
field speed. The supply frequency sets the field speed. So the motor speed is hard to
change without extra electronics. Homes use single-phase AC (one AC supply, not three).
It cannot make a rotating field by itself, so appliances need an extra starting coil.

\subsection{Magnetic braking}
\term{Magnetic braking} is slowing a moving metal object with eddy currents, without
touching it. \Cref{ch5:fig:disc} shows how it works for a spinning disc.

\begin{sidebyside}[0.45]
  \centering
  \begin{tikzpicture}[x=1.0cm, y=1.0cm]
    \filldraw[fill=hscmuted!12, draw=hscink, line width=0.9pt] (0,0) circle (1.9);
    \fill[hscink] (0,0) circle (1.5pt);
    \draw[->, hscmuted, line width=0.9pt] (125:2.15) arc (125:55:2.15);
    \node[font=\footnotesize, text=hscmuted, anchor=south] at (90:2.17) {disc turns clockwise};
    \filldraw[fill=hscfield!14, draw=hscfield, dashed, line width=0.6pt] (0.35,-0.75) rectangle (1.75,0.75);
    \foreach \xx in {0.55,1.55} \foreach \yy in {-0.42,0.42}
      \node[text=hscfield, font=\footnotesize] at (\xx,\yy) {$\times$};
    \node[font=\scriptsize, text=hscfield, anchor=east, align=right] at (0.3,0.55) {$B$ into\\the page};
    \draw[hscseriesB, line width=1.2pt, ->] (1.05,0.75) ++(-90:0.3) arc (-90:250:0.3);
    \draw[hscseriesA, line width=1.2pt, ->] (1.05,-0.75) ++(90:0.3) arc (90:-250:0.3);
    \draw[->, hscink, line width=0.9pt] (0.5,-0.15) -- (1.6,-0.15) node[pos=0.25, above, font=\scriptsize] {$I$};
    \draw[force] (1.05,-0.08) -- (1.05,0.4) node[pos=0.55, right, text=hscforce, font=\footnotesize] {$F$};
    \draw[velocity] (1.85,0.45) -- (1.85,-0.45) node[pos=0.5, right, text=hscvel, font=\footnotesize] {$v$};
    \node[font=\scriptsize, text=hscseriesB, anchor=west] at (1.42,1.28) {metal enters field};
    \node[font=\scriptsize, text=hscseriesA, anchor=west] at (1.42,-1.28) {metal leaves field};
  \end{tikzpicture}
  \captionof{figure}{A metal disc turning through a field into the page ($\times$) on its
    right side.}\label{ch5:fig:disc}
\nextside
  Where the metal \emph{enters} the field, the flux into the page grows. The eddy
  current flows anticlockwise. It makes a field out of the page that opposes the growth.
  Where the metal \emph{leaves}, the flux falls, so the eddy current flows clockwise.
\end{sidebyside}
Inside the field both loops carry current $I$ to the right. This current is in the
field, so it feels a force $F$ that points \emph{up}: opposite to the motion $v$. The
disc slows down.

\textbf{Energy.} The disc's kinetic energy becomes electrical energy of the eddy
currents, then heat in the metal. If the force helped the motion, the disc would speed
itself up and gain energy from nothing. That breaks conservation of energy. So again,
Lenz's law is conservation of energy.

\textbf{Properties of magnetic brakes.}
\begin{itemize}
  \item The braking force grows with speed (faster motion, faster flux change, bigger
        eddy currents). It is strong at high speed and weak at low speed.
  \item At zero speed there is no flux change and no force. So a magnetic brake cannot
        hold an object still. Friction brakes finish the stop.
  \item A stronger field, or thicker metal that conducts better, brakes more strongly.
        Slots cut in the metal break up the eddy-current loops and weaken the braking
        (the same idea as a laminated transformer core).
  \item Uses: roller coasters, drop towers, trains, exercise bikes. Advantages: no
        contact, so no wear; smooth and quiet; works in the wet. Permanent magnets need
        no power, so the brake still works in a power failure.
\end{itemize}

\textbf{Regenerative braking.} Electric cars and trains use \term{regenerative braking}:
the motor acts as a generator, and the battery stores the energy. The wheels turn the
motor's coil, so its induced emf becomes larger than the battery voltage. The current
reverses and charges the battery (kinetic $\to$ electrical $\to$ chemical energy). The
current in the coil feels a force that opposes the rotation (Lenz). So the vehicle
slows, and energy is saved instead of wasted as heat.

\begin{formulacard}[combine={$K = \frac{1}{2}mv^2$}]{Heat made by a magnetic brake}
  \cardformula{Q = \frac{1}{2}mu^2 - \frac{1}{2}mv^2}
  \begin{symbols}
    \sym{$Q$}{heat produced in the metal of the brake}{\si{J}}
    \sym{$m$}{mass of the moving object}{\si{kg}}
    \sym{$u$}{speed before braking}{\si{m.s^{-1}}}
    \sym{$v$}{speed after braking}{\si{m.s^{-1}}}
  \end{symbols}
  \inwords{the heat made by the brake equals the kinetic energy lost.}
  \whentouse{an object slowed by magnetic brakes on a level track, other losses ignored.
    Average heating power: $P = \dfrac{\Delta E}{\Delta t}$ (on the sheet).}
  \whennottouse{on its own when the height changes (a drop tower, a slope). Then add
    the potential energy lost, $\Delta U = mg\Delta h$ (on the sheet).}
  \rearranged{$Q = \dfrac{1}{2}m\left(u^2 - v^2\right)$ \qquad
    $v = \sqrt{u^2 - \dfrac{2Q}{m}}$}
  \tensecond{$m = \qty{2.0}{kg}$, $u = \qty{3.0}{m.s^{-1}}$, $v = \qty{1.0}{m.s^{-1}}$:
    $Q = \frac{1}{2}(2.0)(3.0^2) - \frac{1}{2}(2.0)(1.0^2) = 9.0 - 1.0 = \qty{8.0}{J}$.}
  \derive{use conservation of energy and $K = \frac{1}{2}mv^2$ (on the sheet).
    \begin{align*}
      \tfrac{1}{2}mu^2 &= \tfrac{1}{2}mv^2 + Q && \text{energy before $=$ energy after}\\
      Q &= \tfrac{1}{2}mu^2 - \tfrac{1}{2}mv^2 && \text{subtract $\tfrac{1}{2}mv^2$ from both sides}
    \end{align*}}
  \trap{square each speed before you subtract: $\frac{1}{2}m(u^2 - v^2)$, not
    $\frac{1}{2}m(u - v)^2$. Change \si{km.h^{-1}} to \si{m.s^{-1}} first (divide by 3.6).}
\end{formulacard}

\begin{workedexample}{Braking a roller coaster}\label{ch5:we:coaster}
A \qty{1800}{kg} roller-coaster train enters a level magnetic braking section at
\qty{25}{m.s^{-1}}. It leaves \qty{3.5}{s} later at \qty{8.0}{m.s^{-1}}. Calculate
(a) the heat produced in the brakes and (b) the average braking power.
\given{$m = \qty{1800}{kg}$, \ $u = \qty{25}{m.s^{-1}}$, \ $v = \qty{8.0}{m.s^{-1}}$,
  \ $\Delta t = \qty{3.5}{s}$; level track}
\find{(a) $Q$; \ (b) $P$}
\formula{$Q = \frac{1}{2}mu^2 - \frac{1}{2}mv^2$, \quad $P = \dfrac{\Delta E}{\Delta t}$}
\substitute{\eqlines{\tfrac{1}{2}mu^2 &= \tfrac{1}{2}(\qty{1800}{kg})(\qty{25}{m.s^{-1}})^2 = \qty{562500}{J}\\
  \tfrac{1}{2}mv^2 &= \tfrac{1}{2}(\qty{1800}{kg})(\qty{8.0}{m.s^{-1}})^2 = \qty{57600}{J}\\
  Q &= \qty{562500}{J} - \qty{57600}{J} = \qty{504900}{J}\\
  P &= \frac{\qty{504900}{J}}{\qty{3.5}{s}} = \qty{1.44e5}{W}}}
\answer{(a) \final{$Q = \qty{5.0e5}{J}$} \quad (b) \final{$P = \qty{1.4e5}{W}$}}
\checkstep{$Q$ (\qty{5.0e5}{J}) is less than the starting kinetic energy
  (\qty{5.6e5}{J}) \cmark. $v \div u = 0.32$, so $0.32^2 \approx 10\%$ of the kinetic
  energy is left: about 90\% is removed \cmark. Deceleration
  $(25 - 8.0) \div 3.5 = \qty{4.9}{m.s^{-2}}$, about half of $g$ \cmark.}
\end{workedexample}

\begin{workedexample}{A drop-tower ride}\label{ch5:we:droptower}
A \qty{500}{kg} drop-tower car falls freely, then reaches magnetic brakes. During the
last \qty{20}{m} of its fall, the brakes slow it from \qty{30}{m.s^{-1}} to
\qty{5.0}{m.s^{-1}}. Ignore air resistance and friction. Calculate the heat produced in
the brakes.
\given{$m = \qty{500}{kg}$, \ $u = \qty{30}{m.s^{-1}}$, \ $v = \qty{5.0}{m.s^{-1}}$,
  \ $\Delta h = \qty{20}{m}$, \ $g = \qty{9.8}{m.s^{-2}}$}
\find{the heat $Q$}
\formula{heat $=$ kinetic energy lost $+$ potential energy lost:
  $Q = \left(\frac{1}{2}mu^2 - \frac{1}{2}mv^2\right) + mg\Delta h$}
\substitute{\eqlines{K\text{ lost} &= \tfrac{1}{2}(\qty{500}{kg})(\qty{30}{m.s^{-1}})^2
  - \tfrac{1}{2}(\qty{500}{kg})(\qty{5.0}{m.s^{-1}})^2\\
  &= \qty{225000}{J} - \qty{6250}{J} = \qty{218750}{J}\\
  U\text{ lost} &= \qty{500}{kg} \times \qty{9.8}{m.s^{-2}} \times \qty{20}{m} = \qty{98000}{J}\\
  Q &= \qty{218750}{J} + \qty{98000}{J} = \qty{316750}{J}}}
\answer{\final{$Q = \qty{3.2e5}{J}$}}
\checkstep{$Q$ is bigger than the kinetic energy lost, because the brakes also take the
  energy the car gains by falling \cmark.}
\end{workedexample}

\begin{mistakes}
\begin{itemize}
  \item \xmark\ ``Current is fed to the rotor through brushes.'' \cmark\ Nothing is
        joined to the rotor; its currents are induced.
  \item \xmark\ ``The rotor turns at the same speed as the field.'' \cmark\ It is always
        a little slower (slip). The field really rotates.
  \item \xmark\ ``The magnet attracts the copper or aluminium.'' \cmark\ These metals are
        not magnetic. The force comes from induced eddy currents (Lenz).
  \item \xmark\ ``A magnetic brake can stop a train and hold it still.'' \cmark\ There
        is no braking force at zero speed.
  \item \xmark\ Writing $\frac{1}{2}m(u - v)^2$. \cmark\ $\frac{1}{2}m(u^2 - v^2)$.
\end{itemize}
\end{mistakes}

\begin{examtip}[breakable]
Both topics are often asked as 5--7 mark ``explain'' questions with a diagram. Each
linked step can earn a mark.

\textbf{Induction motor:}
\begin{enumerate}
  \item Parts: stator coils on three-phase AC; a squirrel cage of bars and end rings.
  \item The stator makes a rotating field.
  \item The flux through the loops of bars changes, so currents are induced (Faraday).
  \item The currents are in the field, so the bars feel forces (motor effect).
  \item Lenz: the forces oppose the relative motion, so the rotor follows the field.
  \item Slip: the rotor is always a little slower than the field.
\end{enumerate}
\textbf{Magnetic brake:}
\begin{enumerate}
  \item Draw the metal, the field, the eddy currents and the force opposite to $v$.
  \item The flux through the moving metal changes, so eddy currents are induced
        (Faraday).
  \item Lenz: the force on the eddy currents opposes the motion.
  \item Energy: kinetic $\to$ electrical $\to$ heat.
  \item It cannot hold an object still: no motion, no flux change, no force.
  \item Link to the context, for example: no contact, so no wear.
\end{enumerate}
``Assess'' or ``evaluate'' questions also need a clear judgement at the end.
\end{examtip}

\begin{quickcheck}[breakable]
  \item Explain why an AC induction motor does not need brushes or a commutator.
  \item The field of an induction motor turns at \qty{3000}{rpm}. The rotor turns at
        \qty{2910}{rpm}. (a) Find the slip in rpm and in turns each second. (b) What
        happens to the slip, the induced current and the torque when the load increases?
  \item Explain why magnetic brakes alone cannot hold a train still on a slope.
  \item A \qty{3.0e4}{kg} train carriage is slowed by magnetic brakes from
        \qty{20}{m.s^{-1}} to \qty{12}{m.s^{-1}} on a level track. Calculate the heat
        produced.
  \item A spinning metal disc slows down when a magnet is brought near it. Relate this
        to Lenz's law and the law of conservation of energy.
\end{quickcheck}
%% ------------------------------------------------------------------
\section{Key experiments}\label{ch5:sec:experiments}
\covers{M6.IQ4.01, M6.IQ4.03}

\begin{experiment}{Investigating a simple DC motor}
\aim{To build a simple DC motor, find the job of each part, and see how back emf
  affects the current.}
\apparatus{enamelled copper wire (copper wire with a thin insulating coat), two large
  paper clips, a strong magnet, a \qty{3.0}{V} DC supply, an ammeter, a switch,
  sandpaper, leads, and a small lab DC motor (for the current readings)
  (\cref{ch5:fig:expmotor}).
  \begin{figure}
    \centering
    \begin{tikzpicture}[x=1cm, y=1cm]
      \fill[ground] (-2.6,-0.25) rectangle (2.6,0);
      \draw[hscink, line width=0.6pt] (-2.6,0) -- (2.6,0);
      \filldraw[fill=hscmuted!45, draw=hscink, line width=0.6pt] (-0.45,0) rectangle (0.45,0.28);
      \node[diagramlabel, font=\footnotesize, anchor=south] at (0,0.3) {magnet};
      \draw[hscmuted!80!black, line width=1.6pt] (-1.4,0) -- (-1.4,1.55);
      \draw[hscmuted!80!black, line width=1.6pt] (1.4,0) -- (1.4,1.55);
      \draw[hscmuted!80!black, line width=1.2pt] (-1.4,1.68) circle (0.13);
      \draw[hscmuted!80!black, line width=1.2pt] (1.4,1.68) circle (0.13);
      \draw[hscseriesB!85!black, line width=1.3pt] (0,1.68) ellipse (0.55 and 0.62);
      \draw[hscseriesB!85!black, line width=1.3pt] (0,1.68) ellipse (0.5 and 0.57);
      \draw[hscseriesB!85!black, line width=1.3pt] (-1.62,1.68) -- (-0.55,1.68);
      \draw[hscseriesB!85!black, line width=1.3pt] (0.55,1.68) -- (1.62,1.68);
      \draw[hscink, line width=0.7pt] (-1.4,0.4) -- (-3.2,0.4) -- (-3.2,-0.9) -- (-0.8,-0.9);
      \draw[hscink, line width=0.7pt] (1.4,0.4) -- (3.2,0.4) -- (3.2,-0.9) -- (1.35,-0.9);
      \filldraw[fill=white, draw=hscink, line width=0.7pt] (1.1,-0.9) circle (0.25);
      \node[font=\footnotesize\bfseries] at (1.1,-0.9) {A};
      \draw[hscink, line width=0.7pt] (0.85,-0.9) -- (0.2,-0.9);
      \draw[hscink, line width=1.6pt] (-0.8,-1.04) -- (-0.8,-0.76);
      \draw[hscink, line width=0.9pt] (-0.5,-1.18) -- (-0.5,-0.62);
      \draw[hscink, line width=0.7pt] (-0.5,-0.9) -- (-0.2,-0.9);
      \draw[hscink, line width=0.7pt] (-0.2,-0.9) -- (0.15,-0.72);
      \fill[hscink] (-0.2,-0.9) circle (1.2pt) (0.2,-0.9) circle (1.2pt);
      \node[diagramlabel, font=\footnotesize, anchor=north] at (-0.65,-1.2) {DC supply};
      \node[diagramlabel, font=\footnotesize, anchor=south] at (0.0,-0.68) {switch};
      \node[diagramlabel, font=\footnotesize, anchor=north] at (1.25,-1.2) {ammeter};
      \draw[hscmuted, line width=0.4pt] (-1.0,1.68) -- (-2.4,2.45) node[anchor=south, diagramlabel, font=\footnotesize, align=center] {end fully scraped};
      \draw[hscmuted, line width=0.4pt] (1.0,1.68) -- (2.6,2.45) node[anchor=south, diagramlabel, font=\footnotesize, align=center] {end half scraped\\(acts as a commutator)};
      \draw[hscmuted, line width=0.4pt] (0.1,2.3) -- (0.1,3.0) node[anchor=south, diagramlabel, font=\footnotesize] {coil (20 turns)};
      \draw[hscmuted, line width=0.4pt] (1.4,1.0) -- (2.3,1.15) node[anchor=west, diagramlabel, font=\footnotesize, align=left] {paper-clip support\\(acts as a brush)};
    \end{tikzpicture}
    \caption{A simple DC motor made from a coil of wire, paper clips and a magnet.}
    \label{ch5:fig:expmotor}
  \end{figure}}
\method{\begin{enumerate}[beginpenalty=10000]
  \item Wind about 20 turns into a coil \qtyrange{2}{3}{cm} across, with straight ends as
        the axle. Sand all the insulation off one end, but only the top half off the
        other end. That end switches the current off for the half turn when the torque
        would act the wrong way.
  \item Rest the ends in the supports above the magnet. Connect the supply, switch and
        ammeter in series. Close the switch and push the coil gently. (A single coil can
        stop where the torque is zero, so it may need a start.)
  \item Change one variable at a time (turns, magnets, voltage). Note how easily the coil
        starts and how quickly it gains speed.
  \item Back emf: connect the small lab DC motor in place of the coil. Read the current
        while it runs freely. Then hold its shaft still for 1--2 seconds and read the
        current again.
\end{enumerate}}
\results{the coil spins only when one end is half scraped. More turns, a stronger magnet
  or a larger voltage give more torque, so the coil starts more easily and gains speed
  more quickly. Sample readings for the lab motor: held still \qty{1.5}{A}; running freely
  \qty{0.30}{A}.}
\conclusion{the motor effect gives a torque on the coil. The half-scraped end works like
  a commutator, so the torque always acts the same way. The lab motor draws less current
  when it runs, because its spinning coil makes a back emf that opposes the supply. Held
  still, it has no back emf, so the resistance of the whole circuit is
  $R = \dfrac{3.0}{1.5} = \qty{2.0}{\ohm}$. Running:
  $\varepsilon_\mathrm{back} = V - IR = 3.0 - 0.30 \times 2.0 = \qty{2.4}{V}$.
  (The current in the paper-clip motor also falls when it spins. The main reason is
  different: the half-scraped end switches the current off for half of each turn. Its
  small coil makes a back emf of only about \qty{0.1}{V}.)}
\significance{it shows every part of a DC motor, in a motor you can build in minutes.
  The lab motor makes back emf easy to see.}
\validity{\begin{itemize}[beginpenalty=10000]
  \item \textbf{Validity:} change one variable at a time. Keep the coil size, the magnet
        distance and the voltage the same, and balance the coil.
  \item \textbf{Reliability:} repeat each reading three times and average. Close
        agreement means the results are reliable.
  \item \textbf{Accuracy:} use a digital ammeter and read it when steady. Measure speed
        with slow-motion video or a tachometer (a meter that measures how fast something
        spins), not by eye.
  \item \textbf{Safety:} a motor that is held still draws a large current and gets hot.
        Hold the lab motor still for only 1--2 seconds. Switch off between trials.
\end{itemize}}
\end{experiment}

\textbf{Magnetic braking practical.} The falling-magnet and swinging-plate experiment in
\cref{ch3:chap} is the key practical for magnetic braking. It shows the same physics as the
brakes in \cref{ch5:sec:induction}: eddy currents, Lenz's law, and kinetic or potential
energy changing to heat.

%% ------------------------------------------------------------------
\begin{summarybox}[unbreakable]{\Cref{ch5:chap} in five lines}
\begin{itemize}
  \item A coil in a field feels a torque $\tau = nIAB\sin\theta$, with $\theta$ from the
        normal. It is largest when the plane of the coil is parallel to $B$. A split-ring
        commutator reverses the current every half turn.
  \item A spinning motor coil makes a back emf that opposes the supply:
        $I = \dfrac{V - \varepsilon_\mathrm{back}}{R}$ (not on the sheet). The current is
        largest at switch-on.
  \item A generator coil turned in a field has a changing flux, so an emf is induced.
        Slip rings give AC; a split ring gives DC. Faster turning raises the peak and the
        frequency.
  \item AC induction motor: a rotating field induces currents in the rotor, which follows
        the field a little more slowly (slip).
  \item Magnetic brake: eddy currents oppose the motion and turn kinetic energy into heat.
        There is no braking force at zero speed.
\end{itemize}
\end{summarybox}
